% concatenated from main_18.03.26.tex

\documentclass[11pt,leqno]{article}
\usepackage{amsmath, amscd, amsthm, amssymb, graphics, xypic, mathrsfs, setspace, fancyhdr, times, bm, enumitem}
%\xyoption{all}
\usepackage[letterpaper,top=1.05in, bottom=1.05in, left=1.05in, right=1.05in]{geometry}
\usepackage[colorlinks=true,pagebackref=true]{hyperref} % new hyperref package added by Jean
\hypersetup{backref}

\newcommand{\william}[1]{\textcolor{red}{#1}} %notes by William
\newcommand{\jean}[1]{\textcolor{blue}{#1}} %notes by Jean

%\include{commands}

\newcommand{\setof}[1]{\{ #1 \}}
\newcommand{\tensor}{\otimes}
\newcommand{\colim}{\operatorname{colim}}
\newcommand{\Spec}{\operatorname{Spec}}
\newcommand{\Aut}{{\mathbf{Aut}}}
\newcommand{\isomto}{{\stackrel{\sim}{\;\longrightarrow\;}}}
\newcommand{\isomt}{{\stackrel{{\scriptscriptstyle{\sim}}}{\;\rightarrow\;}}}
\newcommand{\smallsim}{{\scriptscriptstyle{\sim}}}
\newcommand{\sma}{{\scriptstyle{\wedge}}\;}
\newcommand{\coker}{\operatorname{coker}}
\newcommand{\Singaone}{\operatorname{Sing}^{\aone}\!\!}

\renewcommand{\O}{{\mathcal O}}
\renewcommand{\hom}{\operatorname{Hom}}
\newcommand{\EH}{\operatorname{EH}}
\newcommand{\Xa}{{\mathcal{X}_{\underline{a}}}}
\newcommand{\real}{{\mathbb R}}
\newcommand{\cplx}{{\mathbb C}}
\newcommand{\Q}{{\mathbb Q}}
\newcommand{\Z}{{\mathbb Z}}
\newcommand{\N}{{\mathbb N}}
\newcommand{\A}{{\mathbb A}}
\newcommand{\Proj}{{\mathbb P}}
\newcommand{\Jo}{{\mathbb J}}
\newcommand{\aone}{{\mathbb A}^1}
\newcommand{\pone}{{\mathbb P}^1}
\newcommand{\GL}{{\mathbf{GL}}}
\newcommand{\PGL}{{\mathbf{PGL}}}
\newcommand{\SL}{{\mathbf{SL}}}
\newcommand{\Stilde}{{\tilde{{\mathcal S}}}}
\renewcommand{\1}{{\rm 1\hspace*{-0.4ex}%
\rule{0.1ex}{1.52ex}\hspace*{0.2ex}}}

\newcommand{\ga}{{{\mathbb G}_{a}}}
\newcommand{\gm}{{{\mathbb G}_{m}}}
\renewcommand{\L}{{\mathcal L}}
\newcommand{\MW}{\mathrm{MW}}
\newcommand{\So}{{\mathbb S}^0}
\newcommand{\et}{\text{\'et}}
\newcommand{\ho}[1]{\mathscr{H}({#1})}
\newcommand{\hop}[1]{\mathscr{H}_{\bullet}({#1})}
\newcommand{\DM}{{\mathbf{DM}}}
\newcommand{\M}{{\mathbf{M}}}
\newcommand{\Sp}{{\mathscr{Sp}}}
\newcommand{\bpi}{\bm{\pi}}
\newcommand{\bmu}{\bm{\mu}}
\newcommand{\piaone}{{\bpi}^{\aone}}
\newcommand{\Daone}[1]{{\mathrm D}_{\aone}(#1)}
\newcommand{\eff}{\operatorname{eff}}
\newcommand{\Nis}{\operatorname{Nis}}
\newcommand{\Zar}{\operatorname{Zar}} % Zariski
\newcommand{\SH}{{\mathbf{SH}}}
\newcommand{\CH}{{\widetilde{CH}}}
\newcommand{\Ch}{{\mathscr{Ch}}}
\newcommand{\Fun}{\mathscr{Fun}}
\renewcommand{\deg}{\operatorname{deg}}
\newcommand{\tdeg}{\widetilde{\deg}}
\newcommand{\Laone}{\operatorname{L}_{\aone}}

\newcommand{\Shv}{{\mathscr{Shv}}}
\newcommand{\Sm}{\mathrm{Sm}}
\newcommand{\Cor}{\mathscr{Cor}}
\newcommand{\Spc}{\mathrm{Spc}}
\newcommand{\Mod}{\mathscr{Mod}}
\newcommand{\Ab}{\mathscr{Ab}}
\newcommand{\Gr}{\mathscr{Gr}}
\newcommand{\Set}{\mathscr{Set}}
\newcommand{\K}{{{\mathbf K}}}
\newcommand{\KMW}[1]{{{\mathbf K}_{#1}^{\MW}}}
\newcommand{\KM}[1]{{{\mathbf K}_{#1}^{\mathrm{M}}}}
\newcommand{\Pic}{\operatorname{Pic}}
\newcommand{\Faone}{{\mathbf{F}}_{\aone}}
\renewcommand{\H}{{{\mathbf H}}}
\newcommand{\Alg}{{\mathbf{Alg}}}
\newcommand{\Sym}{{\operatorname{Sym}}}
\newcommand{\<}[1]{{\langle}#1 {\rangle}}
\newcommand{\simpspc}{\Delta^{\circ}\Spc}
\newcommand{\limdir}{\underrightarrow {\operatorname{\mathstrut lim}}}
\newcommand{\hsnis}{\mathscr{H}_s^{\Nis}(k)}
\newcommand{\hspnis}{\mathscr{H}_{s,\bullet}^{\Nis}(k)}
\newcommand{\hspet}{\mathscr{H}_{s,\bullet}^{\et}(k)}
\newcommand{\hsptau}{\mathscr{H}_{s,\bullet}^{\tau}(k)}
\newcommand{\hset}{\mathscr{H}_s^{\et}(k)}
\newcommand{\hstau}{\mathscr{H}_s^{\tau}(k)}
\newcommand{\het}[1]{\mathscr{H}^{\et}(#1)}
\newcommand{\hpet}[1]{\mathscr{H}^{\et}_{\bullet}(#1)}
\newcommand{\dkset}{{\mathfrak D}_k-{\mathcal Set}}
\newcommand{\F}{{\mathcal F}}
\newcommand{\fkset}{{\mathcal F}_k-{\mathcal Set}}
\newcommand{\fkrset}{{\mathcal F}_k^r-{\mathcal Set}}
\newcommand{\Zn}{\Z \langle n \rangle}

\newcommand{\simpnis}{{\Delta}^{\circ}\Shv_{\Nis}({\mathcal Sm}_k)}
\newcommand{\simpmod}{{\Delta}^{\circ}\Mod}
\newcommand{\simpet}{{\Delta}^{\circ}\Shv_{\text{\'et}}({\mathcal Sm}_k)}
\newcommand{\simptau}{{\Delta}^{\circ}\Shv_{\tau}({\mathcal Sm}_k)}

\newcommand{\comment}[1]{\marginpar{\begin{tiny}{#1}\end{tiny}}}
\newcommand{\hooklongrightarrow}{\lhook\joinrel\longrightarrow}
\newcommand{\twoheadlongrightarrow}{\relbar\joinrel\twoheadrightarrow}


\newcommand{\kbar}{{\overline{k}}}

\newcommand{\Addresses}{{% additional braces for segregating \footnotesize
 \bigskip
 \footnotesize

 J.~Fasel, Institut Fourier, UMR 5582, Université Grenoble Alpes, CS 40700, 38058 Grenoble cedex 9, France \textit{E-mail address:} \url{jean.fasel@univ-grenoble-alpes.fr}

 \medskip

 W.P.~Hornslien, Institut Fourier, UMR 5582, Université Grenoble Alpes, CS 40700, 38058 Grenoble cedex 9, France \textit{E-mail address:} \url{w.hornslien.math@proton.me}
}}

\newcounter{intro}
\setcounter{intro}{1}

%\renewcommand{\baselinestretch}{1.5}
\theoremstyle{plain}
\newtheorem{thm}{Theorem}[subsection]
\newtheorem{scholium}[thm]{Scholium}
\newtheorem{lem}[thm]{Lemma}
\newtheorem{cor}[thm]{Corollary}
\newtheorem{prop}[thm]{Proposition}
\newtheorem*{claim*}{Claim}  %Claim
\newtheorem{claim}[thm]{Claim}
\newtheorem{question}[thm]{Question}
\newtheorem{problem}[thm]{Problem}
\newtheorem{answer}[thm]{Answer}
\newtheorem{conj}[thm]{Conjecture}
\newtheorem*{thm*}{Theorem}
\newtheorem*{problem*}{Problem}

\newtheorem{thmintro}{Theorem}
\newtheorem{propintro}[thmintro]{Proposition}
\newtheorem{corintro}[thmintro]{Corollary}
\newtheorem{problemintro}[thmintro]{Problem}
\newtheorem{questionintro}[thmintro]{Question}
\newtheorem{defnintro}[thmintro]{Definition}
\newtheorem{scholiumintro}[thmintro]{Scholium}
\newtheorem{lemintro}[thmintro]{Lemma}
\newtheorem{conjintro}[thmintro]{Conjecture}

\theoremstyle{definition}
\newtheorem{defn}[thm]{Definition}
\newtheorem{construction}[thm]{Construction}
\newtheorem{notation}[thm]{Notation}

\theoremstyle{remark}
\newtheorem{rem}[thm]{Remark}
\newtheorem{remintro}[thmintro]{Remark}
\newtheorem{extension}[thm]{Extension}
\newtheorem{ex}[thm]{Example}
\newtheorem{entry}[thm]{}

\numberwithin{equation}{section}

\begin{document}
\pagestyle{fancy}
\renewcommand{\sectionmark}[1]{\markright{\thesection\ #1}}
\fancyhead{}
\fancyhead[LO,R]{\bfseries\footnotesize\thepage}
\fancyhead[LE]{\bfseries\footnotesize\rightmark}
\fancyhead[RO]{\bfseries\footnotesize\rightmark}
\chead[]{}
\cfoot[]{}
\setlength{\headheight}{1cm}

\author{ Jean Fasel\thanks{Jean Fasel was partially supported by the project IRGA2024 –MOTOR – G7H-IRG24E97.}\and William Hornslien\thanks{William Hornslien was partially supported by the project IRGA2024 –MOTOR – G7H-IRG24E97.} }

\title{{\bf Exotic Hopf maps, weight shifting and applications to vector bundles}}
\date{}
\maketitle

\begin{abstract}
Using motivic homotopy theory we produce several explicit polynomial representatives of the suspension of the Hopf map defined over $\Z$. We derive from this computation an explicit rank $2$ vector bundle on the Jouanolou device of $\mathbb{P}^3_\Z$.
\end{abstract}

\begin{footnotesize}
\setcounter{tocdepth}{1}
\tableofcontents
\end{footnotesize}

\section{Introduction}


This paper can be viewed as an explicit version of some results in \cite{Asok20b}, where the problem of producing complex polynomial representatives of homotopy groups of spheres was considered. Recall that P. Baum proved in \cite{Baum67} that any element of the \emph{stable} homotopy groups of spheres can be represented by an admissible quadratic map, i.e., a real polynomial map of degree at most $2$. However, Wood later observed \cite{Wood68} that there are no non-constant polynomial maps $S^n\to S^m$ for $n\geq 2m$, showing that Baum's result fails in the unstable range. This led to the study of complex polynomial representatives, using the fact that the complex algebraic sphere $S_{\mathbb{C}}^n$, viewed as a manifold, is homotopy equivalent to the real sphere $S^n$. At present, no element of a homotopy group of spheres is known that cannot be represented by such a complex polynomial map.

In \cite[Proposition 2.2.1]{Asok20b}, it was shown that the suspension of the Hopf map $S^4\to S^3$ and the square of the Hopf map $S^5\to S^3$ admit complex polynomial representatives. The method relied on motivic homotopy theory and a weight-shifting procedure that we briefly recall. If $k$ is a field and $n\geq 1$, consider the smooth quadrics $Q_{2n-1}:=\{\sum x_iy_i=1\}\subset \A^{2n}_k$ (with coordinates $x_1,\ldots,x_n,y_1,\ldots,y_n$) and $Q_{2n}:=\{\sum x_iy_i=z(1-z)\}\subset \A^{2n+1}_k$ (with coordinates $x_1,\ldots,x_n,y_1,\ldots,y_n,z$). It is well known (e.g. \cite{Asok14}) that these quadrics model motivic spheres in the (unstable) motivic homotopy category $\mathcal{H}(k)$ of Morel–Voevodsky \cite{Morel99}. Moreover, for any smooth affine $k$-scheme $X$, the map $\mathrm{Hom}_{\mathrm{Sch}_k}(X,Q_i)\to [X,Q_i]_{\A^1}:=\mathrm{Hom}_{\mathcal{H}(k)}(X,Q_i)$ is surjective for $i\geq 1$ (\cite[Theorem 4.2.1]{Asok15b}, \cite[Corollary 3.1.1]{Asok22}).

If $k\subset \cplx$, the complex realization functor \cite[\S 3.3]{Morel99}
\[
r_{\cplx}\colon \mathcal{H}(\cplx)\to \mathcal{H}
\]
(left Kan extended from $X\mapsto X(\cplx)$ endowed with the complex topology) sends $Q_i$ to $S^i$. Hence, we obtain a composite map
\[
\mathrm{Hom}_{\mathrm{Sch}_k}(Q_j,Q_i)\to [Q_j,Q_i]_{\A^1}\to \pi_j(S^i)
\]
for any $i,j\geq 1$. Thus, to show that $\alpha\in \pi_j(S^i)$ admits a polynomial representative, it suffices to produce a lift $\alpha'\in [Q_j,Q_i]_{\A^1}$. Since $Q_j$ is a motivic sphere, we have identifications with bi-graded homotopy groups
\[
[Q_j,Q_i]_{\A^1}=\begin{cases} \piaone_{n,n}(Q_i)(k) & \text{ if $j=2n$, }  \\ \piaone_{n-1,n}(Q_i)(k) & \text{ if $j=2n-1$. }\end{cases}
\]
These groups can sometimes be computed using fiber sequences analogous to those in topology. For example, setting $i=3,j=4$ yields
\[
\piaone_{2,2}(Q_3)=\piaone_{2,2}(\mathrm{SL}_2)=\piaone_{2,2}(\A^2\smallsetminus\{0\}),
\]
which fits in an exact sequence (\cite[Theorem 3.3, Proposition 4.4]{Asok12a})
\[
\KM 2(k)/12\times_{\KM 2(k)/2}\mathbf{I}^2(k)\to \piaone_{2,2}(Q_3)(k)\to \mathbf{GW}_1^0(k)\to 0.
\]
If $k$ has characteristic $\neq 2$, then $\mathbf{GW}_1^0(k)=\mathrm{KO}_1(k)\simeq \Z/2\times k^\times/(k^{\times})^2$. When $k=\cplx$, we obtain in particular $\piaone_{2,2}(Q_3)(\cplx)\simeq \Z/2$, and the generator maps to the generator of $\pi_4(S^3)=\Z/2$, namely the suspension of the Hopf map.

In \cite{Asok20b}, a different approach called \emph{weight shifting} was used to produce a motivic lift of the suspension of the Hopf map. We observe that $\piaone_{1,3}(Q_3)=\mathbf{W}$ (the Witt sheaf) and that its generator, as a $\KMW 0$-module, is $2_{\epsilon}$-torsion (with $2_{\epsilon}=\langle 1,-1\rangle\in \mathrm{GW}(k)$). It is shown in \cite[Theorem 2.1.11]{Asok20b} that choosing $-1$ as a primitive $2$nd root of unity induces a morphism of sheaves
\[
\mathbf{W}={}_{2_\epsilon}\piaone_{1,3}(Q_3)\to \piaone_{2,2}(Q_3)/[-1]\piaone_{2,3}(Q_3).
\]
By construction, the image of the generator realizes to the suspension of the Hopf map. Although conceptually natural, this construction is difficult to make explicit and a geometric realization was missing.

In this article we give two explicit constructions of morphisms of quadrics $Q_4\to Q_3$, called \emph{exotic Hopf maps}, whose complex realization is the suspension of the Hopf map. The first construction uses the computation of $\piaone_{2,2}(Q_3)$ above, specifically the morphism $\piaone_{2,2}(Q_3)\to \mathbf{GW}_1^0(k)$ and the summand $\Z/2\subset \mathrm{GW}_1^0(k)$. We first note that $\mathrm{GW}_1^0(k)=\piaone_{2,2}(\mathrm{Sp})$, where $\mathrm{Sp}$ denotes the infinite symplectic group defined over $\Z$. Hence
\[
\piaone_{2,2}(\mathrm{Sp})=[Q_4,\mathrm{Sp}]_{\aone}=\mathrm{GW}_1^2(Q_4)=\mathrm{KSp}_1(Q_4).
\]
Using the edge map in the coniveau spectral sequence for Hermitian $K$-theory, we construct an explicit symplectic matrix $M\in \mathrm{Sp}_8(Q_4)$ representing the generator of $\Z/2$. Connectivity arguments show that this matrix can be reduced by elementary symplectic operations, although carrying this out explicitly is somewhat delicate. Nevertheless, we obtain the following theorem.

\begin{thmintro}
The exotic Hopf map $Q_4\to Q_3$ defined (over $\Z$) by the matrix $\begin{pmatrix}
a & b \\ c & d
\end{pmatrix}$ with

\begin{align*}
a =& 16x_2^3y_1^2y_2+16x_2^2y_1^2z-16x_2^2y_1y_2z-16x_2^2y_1^2-8x_2^2y_1y_2-16x_2^2y_2^2-24x_2y_1z^2 \\ &+20x_2y_1z-8x_2y_2z+8z^3+4x_2y_1-8x_1y_2+12x_2y_2-8z^2-2z+1, \\
b =& -32x_2^3y_2z-32x_1x_2^2y_2+32x_2^3y_2-32x_2^2z^2-48x_1x_2z+64x_2^2z-16x_1^2+40x_1x_2-32x_2^2, \\
c =& -2x_2y_1^3-4x_2y_1^2y_2+2y_1^2z+4y_1y_2z+y_1^2+2y_1y_2+4y_2^2, \\
d =& 4x_2y_1z+8x_2y_2z-4x_2y_1+8x_1y_2-12x_2y_2-4z^2+2z+1.
\end{align*}
complex realizes to the suspension of the Hopf map.
\end{thmintro}

The second construction uses the weight-shifting procedure. Starting from the universal rank $2$ bundle on $Q_5\simeq \A^3\smallsetminus{0}$ classified by a map $Q_5\to \mathrm{BSL}_2$, adjunction yields a map $S^{1,3}\to \mathrm{SL}_2=Q_3$. Since $S^{1,3}$ has no convenient model as a smooth scheme, we instead use a presentation as a quotient of $Q_3\times \gm{}$ (as Nisnevich sheaves of sets). This leads to a morphism of schemes
\[
T\colon Q_3\times \gm{}\to \mathrm{SL}_2=Q_3
\]
which is trivial on $Q_3\vee \gm{}$. Considering a model of the multiplication-by-$2$ map $m'\colon Q_3\to Q_3$, one checks that the composite
\[
Q_3\times \gm{}\xrightarrow{m'\times 1}Q_3\times \gm{}\to Q_3
\]
is homotopically trivial. This produces a morphism
\[
\phi\colon Q_3\times \gm{}\times \A^1\to Q_3
\]
with $\phi(0)=\mathrm{Id}$ and $\phi(1)=T\circ m'$. Evaluating at the point $-1\in\gm{}$ gives a morphism $G\colon Q_3\times\A^1\to Q_3$ satisfying $G(0)=G(1)=\mathrm{Id}$ and then a map $(Q_3\times \A^1)/(Q_3\times S^0)\to Q_3$. The source formally resembles a suspension of $Q_3$, hence $Q_4$. After suitable modifications and using a morphism $Q_3\times(\A^1\smallsetminus\{0,1\})\to Q_4$ constructed in Section \ref{sec:basicsquadrics}, we obtain the desired morphism $Q_4\to Q_3$. The resulting map has relatively high degree, as illustrated by the following theorem.

\begin{thmintro}
The exotic Hopf map $Q_4\to Q_3$ defined (over $\Z$) by the unimodular row $(a',b')$ with
\begin{align*}
a' =& 144x_2^2z^9+144x_2y_1z^9+72x_1x_2z^8-384x_2^2z^8-240x_2y_1z^8-72x_2 y_2z^8-72z^{10}-192x_1x_2z^7\\
&+112x_2^2z^7-128x_2y_1z^7+120x_2y_2z^7+192 z^9+92x_1x_2z^6+384x_2^2z^6+256x_2y_1z^6+28x_2y_2z^6\\
&-92z^8+96x_1x_2z ^5-208x_2^2z^5+48x_2y_1z^5-24y_1^2z^5-68x_2y_2z^5-96z^7-40x_1x_2z^4\\
&- 64x_2^2z^4-16x_2y_1z^4-16y_1^2z^4-28x_2y_2z^4+12y_1y_2z^4+28z^6-32 x_1x_2z^3\\
&-48x_2^2z^3-64x_2y_1z^3+16y_1^2z^3+4x_2y_2z^3+8y_1y_2z^3+48 z^5-12x_1x_2z^2\\
&+64x_2^2z^2+16y_1^2z^2+16x_2y_2z^2-2y_1y_2z^2+18z^4+16x_1x_2z+8y_1^2z\\
&-4y_1y_2z-24z^3-2y_1y_2-2z^2+1, \\
b' =& -144x_1x_2z^9+144x_2y_2z^9-72x_1^2z^8+384x_1x_2z^8+72x_1y_2z^8-240x_2 y_2z^8+192x_1^2z^7\\
&- 112x_1x_2z^7-120x_1y_2z^7-128x_2y_2z^7-92x_1^2z^6 -384x_1x_2z^6-28x_1y_2z^6\\
&+256x_2y_2z^6-96x_1^2z^5+208x_1x_2z^5+68x_1 y_2z^5+48x_2y_2z^5-24y_1y_2z^5\\
&-24z^7+40x_1^2z^4+64x_1x_2z^4+28x_1y_2 z^4-16x_2y_2z^4-16y_1y_2z^4\\
&+12y_2^2z^4+32z^6+32x_1^2z^3+48x_1x_2z^3- 4x_1y_2z^3-64x_2y_2z^3\\
&+16y_1y_2z^3+8y_2^2z^3+24z^5+12x_1^2z^2-64x_1 x_2z^2-16x_1y_2z^2\\
&+16y_1y_2z^2-2y_2^2z^2-32z^4-16x_1^2z+8y_1y_2z-4 y_2^2z-4z^3-2y_2^2+4z.
\end{align*}
complex realizes to the suspension of the Hopf map.
\end{thmintro}

At present we do not know whether the maps in Theorems 1 and 2 are homotopic over $\Z$. Both reduce to the identity over $\mathbb{F}_2$, since the generator of $\mathrm{GW}_1^0(k)$ in the first construction and $-1$ in the second become trivial modulo $2$. Finally, we note that Turiel constructs in \cite{Turiel07} a suspension procedure for certain maps between the algebraic spheres $\mathbb{S}_k^n=\Spec(k[x_1,\ldots,x_{n+1}]/(\Sigma x_i^2-1))$, including the Hopf map. Since $k[\mathbb{S}^{n}_\cplx]\cong k[Q_n]$, this method also gives a map $Q_4\to Q_3$, defined over $\Z[i]$, whose complex realization is the suspension of the Hopf map. The resulting unimodular row $(a^{\prime\prime},b^{\prime\prime})$ is

\begin{align*}
a^{\prime\prime} =& -2x_1^2z^2i-4x_1x_2z^2i-2x_2^2z^2i+2y_1^2z^2i-4y_1y_2z^2i+2y_2^2z^2i-4x_2y_1z^2+4x_1y_2z^2\\
&+8x_2y_2z^2+4z^4+2x_1^2zi+4x_1x_2zi+2x_2^2zi-2y_1^2zi+4y_1y_2zi-2y_2^2zi+4x_2y_1z\\
&-4x_1y_2z-8x_2y_2z-8z^3+x_1^2i+4x_1x_2i+x_2^2i-y_1^2i+4y_1y_2i-y_2^2i\\
&+2x_2y_1-2x_1y_2-8x_2y_2+4z,\\
b^{\prime\prime} =& 2x_1^2z^2i+4x_1x_2z^2i+2x_2^2z^2i+2y_1^2z^2i-4y_1y_2z^2i+2y_2^2z^2i-2x_1^2zi-4x_1x_2zi\\
&-2x_2^2zi-2y_1^2zi+4y_1y_2zi-2y_2^2zi+4z^3i-x_1^2i-4x_1x_2i-x_2^2i\\
&-y_1^2i+4y_1y_2i-y_2^2i-6z^2i+i.
\end{align*}
In Section \ref{sec:pone-suspension}, we develop a method for $\pone$-suspending maps between quadrics. Given a map $f \in \hom_{\mathrm{Sch}_k}(Q_i, Q_j)$, we produce a new map in $\hom_{\mathrm{Sch}_k}(Q_{i+2}, Q_{j+2})$ which has the homotopy class of $\Sigma_{\pone}f$. We use this to give a procedure for constructing a map representing any homotopy class in $[Q_i,Q_i]_{\A^1} \cong \mathrm{GW}(k)$ for $i \geq3$.


In the second part of the paper, we apply our construction of the exotic Hopf map to obtain an explicit rank $2$ vector bundle on a convenient replacement of $\mathbb{P}^3$. For context, recall that Atiyah and Rees observed in \cite{Atiyah76} that the bundle coming from the composite 
\[
\cplx\Proj^3 \to S^6 \xrightarrow{\Sigma^3 \eta }  S^5 \to BU(2)
\]
had trivial Chern classes, but was indecomposable. They then showed that all topological bundles on $\cplx \Proj^2$ are uniquely determined by their Chern classes and a $\Z_2$-invariant they called the $\alpha$-invariant. Horrocks \cite{Horrocks68} came up with a way of constructing rank $2$ bundles on $\cplx\Proj^3$ and Atiyah and Rees proved that all complex rank $2$ topological vector bundles on $\cplx\Proj^3$ can be obtained by starting with a sum of line bundles and then iteratively tensoring with a line bundle and applying Horrocks' construction \cite[Theorem 1.1]{Atiyah76}. In \cite{Asok20b}, it was shown that the same story holds in (unstable) motivic homotopy theory. Namely, there is a composite
\[
\Proj^3_k \to (\pone_k)^{\wedge 3} \xrightarrow{\Sigma_{\Proj^1}\eta' }  \A^3_k\smallsetminus\{0\} \to \mathrm{BSL}_2
\]
where the first map is the ``clutching'' map, the second one is the class of any of the exotic Hopf maps constructed above and the third one is the map classifying the universal rank $2$ bundle on $ \A^3_k\smallsetminus\{0\}$. However, by construction this composite yields a class in $[\Proj^3_k,\mathrm{BSL}_2]_{\A^1}$, i.e., a motivic vector bundle on $\mathbb{P}^3_k$. Such a bundle can be seen as an actual algebraic vector bundle on a convenient affine replacement of $\mathbb{P}^3_k$, the so-called \emph{Jouanolou device} of $\mathbb{P}^3_k$ denoted by $\mathbb{J}^3_k$ below. Up to homotopy, the above composite then reads as
\[
\mathbb{J}^3_k \to Q_6 \xrightarrow{\Sigma_{\Proj^1}\eta' }  Q_5 \to \mathrm{BSL}_2
\]
where the first map is an explicit model of the clutching map provided in Section \ref{sec:collapsemap} below, and the results in the first part of the paper allow us to obtain an explicit idempotent $3\times 3$-matrix of rank $2$ on $\mathbb{J}^3_k$ that is the algebraic construction of the topological composite above.  If $k=\cplx$, we prove that it has the expected algebraic invariants: trivial Chern classes and non-trivial Atiyah-Rees invariant. In view of Horrock's construction, the vector bundle on $\mathbb{J}^3_k$ is pulled-back from $\Proj^3_k$ in case $k=\cplx$, but we don't know if this is the case over an arbitrary field. 

%\jean{Yes, I don't think that we have to add any more detailed description. Also, it would be good to mention explicitly the descent problem. After all, we don't know if the bundles constructed descend to $\Proj^3$ over any field.}
%\william{This introduction is written in a very chronological way, thus we might not need a paragraph saying what happens in each section? I haven't made any big changes, just some small words here and there. I also got a job offer today, so my unemployed era ends soon.}
%\william{Check out this paper by Golasiński and Ruiz: \href{https://www.researchgate.net/publication/225603144_Polynomial_and_Regular_Maps_into_Grassmannians}{Polynomial and regular maps into grassmannians}, they're doing some very affine motivic geometric things. I think we might be able to use the addition structures on $[Q_2,Q_2]$ and $[Q_3,Q_3]$ paired with $\Proj^1$-suspension to give som nice representatives of complex polynomial maps between spheres with certain Brouwer degrees. I just have to browse the literature enough to understand what is actually new.}
%\jean{Congratulations for the job! Yes, it would be great to write down explicit morphisms with given Brouwer degree. I'll have a look at the paper.} \william{Thanks! I'm still up for a meeting this Thursday.}
%Atiyah and Rees observed that the bundle coming from the composite 
%\[
%\cplx\Proj^3 \to S^6 \xrightarrow{\Sigma^3 \eta }  S^5 \to BU(2)
%\]
%had trivial Chern classes, but was indecomposable. They then showed that all (algebraic?) bundles on $\cplx \Proj^2$ are uniquely determined by their Chern classes and a $\Z_2$-invariant they called the $\alpha$-invariant. Horrocks \cite{Horrocks68} came up with a way of constructing rank $2$ bundles on $\cplx\Proj^3$ and Atiyah and Rees proved that all complex rank $2$ topological vector bundles on $\cplx\Proj^3$ can be obtained by starting with a sum of line bundles and then iteratively tensoring with a line bundle and applying Horrocks' construction \cite[Theorem 1.1]{Atiyah76}. 


%The middle set is by definition the motivic homotopy group $\piaone_{n}(Q_i)_{-n}(\cplx)$ if $j=2n$ is even, and $\piaone_{n-1}(Q_i)_{-n}(\cplx)$ if $j=2n-1$ is odd, and it

%The existence of real algebraic representatives for elements of homotopy groups of spheres was studied in the stable setting by [Bau67] and, independently, by R. Wood. Indeed, all elements in the stable ho- motopy groups of spheres admit polynomial representatives (see Baum [Bau67, Corollary 2.11]; Wood’s results remain unpublished). The situation regarding existence of real algebraic representatives of unstable homotopy classes is different. Indeed, R. Wood [Woo68] showed that there are no non-constant polyno- mial maps Sn → Sm when n ≥ 2m. In particular, this means that as soon as p is a prime number ≥ 3, the classes α1 and α12 cannot be represented by morphisms of real algebraic spheres. In fact, even when p = 2, Wood shows that α1 (in this case a generator of π4(S3) ∼= Z/2) admits no real algebraic representative [Woo68, Theorem 2].
%Given the paucity of real algebraic representatives of homotopy classes, Wood later broadened the scope of the representation problem [Woo93] by studying “complex” polynomial representatives. Indeed, if we consider Sn as a complex algebraic variety, then the set of complex points of Sn is a complex manifold diffeomorphic to the tangent bundle of the standard (smooth) n-sphere (in particular, homotopy equivalent to the standard (smooth) n-sphere). By a complex algebraic representative of a homotopy class α ∈ πn(Sn), we mean a morphism of complex algebraic varieties Sn → Sm such that the associated continuous map obtained by taking complex points is in the homotopy class α. Wood writes [Woo93, Question 2]: “What is the first element in the homotopy groups of spheres which cannot be represented by a polynomial map of quadrics?” and then remarks that one cannot rule out the possibility that all elements of the unstable homotopy groups of spheres admits complex polynomial representatives. We present ev- idence in support of this latter possibility, so that the situation regarding existence of complex algebraic representatives of homotopy classes appears to be essentially the opposite of that for real algebraic repre- sentatives.

%Proposition 2.2.1. If k = C, then there exist morphisms of quadrics Q4 → Q3 and Q5 → Q3 providing polynomial representatives of the suspensions of the Hopf map and of the square of the Hopf map.
%Proof. Considerthemapη:A2∖0−→P1.TheG -suspensionofηyieldsamapΣ η:S1∧ G∧3 −→ m Gmm
%A2 ∖ 0. Combining [Mor12, Theorems 6.13 and 6.40], we know that πA1 (A2 ∖ 0)(k) ∼= W(k). In ∼ 1,3
%particular, if k is algebraically closed, then W(k) = Z/2 and we conclude that ΣGm η is a 2-torsion class.
%Takingk=CandappealingtoTheorem2.1.12thereisanelementτΣGmη:P1∧2 →A2∖0.Appealto Lemma 2.1.13 shows that this element is mapped under complex realization to the suspension of η.
%Now, by [ADF17, Proposition 2.1.2] we know that Q4 ∼= P1∧ 2 in H (k). Likewise, Q3 ∼= A2 ∖ 0 in H (k). Then, by [AHW18, Theorem 4.2.1], since the map π0(SingA1Q3)(R) → [Spec R, Q3]A1 is a bijection for any smooth affine k-algebra R, we conclude that τΣGmη is represented by a morphism Q4 → Q3 as required.
%For the second statement, observe that composing ΣP1 η and τ ΣGm η yields an element of [A3 ∖ 0, A2 ∖ 0]A1 that is mapped by complex realization to the suspension of the Hopf map. Again, by [AHW18, Theorem 4.2.1], this element is represented by a morphism of quadrics Q5 → Q3.

%\subsection{An aside on algebraic maps between spheres}
%\label{section:algebraicSphereMaps}
%\jean{I think that this discussion would be nicer in the introduction. Also, we need precise references.}
%Turiel \william{ref} produces a way of suspending maps between the algebraic spheres $\mathbb{S}_k^n = \Spec(k[x_1,\ldots,x_{n+1}]/(\Sigma x_i^2 -1)$.
%
%There's an isomorphism of rings $k[\mathbb{S}^{n}_\cplx]\cong k[Q_{n}]$ and we can use Turiel's method to produce a map $Q_4 \to Q_3$ which complex realizes to the suspension of the Hopf map. 
%
%% -2x_1^2z^2i-4x_1x_2z^2i-2x_2^2z^2i+2y_1^2z^2i-4y_1y_2z^2i+2y_2^2z^2i-4x_2y_1z^2+4x_1y_2z^2+8x_2y_2z^2+4z^4+2x_1^2zi+4x_1x_2zi+2x_2^2zi-2y_1^2zi+4y_1y_2zi-2y_2^2zi+4x_2y_1z-4x_1y_2z-8x_2y_2z-8z^3+x_1^2i+4x_1x_2i+x_2^2i-y_1^2i+4y_1y_2i-y_2^2i+2x_2y_1-2x_1y_2-8x_2y_2+4z&2x_1^2z^2i+4x_1x_2z^2i+2x_2^2z^2i+2y_1^2z^2i-4y_1y_2z^2i+2y_2^2z^2i-2x_1^2zi-4x_1x_2zi-2x_2^2zi-2y_1^2zi+4y_1y_2zi-2y_2^2zi+4z^3i-x_1^2i-4x_1x_2i-x_2^2i-y_1^2i+4y_1y_2i-y_2^2i-6z^2i+i
%The map is given by the unimodular row $(a,b)$, where  
%\begin{align*}
%    a =& -2x_1^2z^2i-4x_1x_2z^2i-2x_2^2z^2i+2y_1^2z^2i-4y_1y_2z^2i+2y_2^2z^2i-4x_2y_1z^2+4x_1y_2z^2\\&+8x_2y_2z^2+4z^4+2x_1^2zi+4x_1x_2zi+2x_2^2zi-2y_1^2zi+4y_1y_2zi-2y_2^2zi+4x_2y_1z-4x_1y_2z\\&-8x_2y_2z-8z^3+x_1^2i+4x_1x_2i+x_2^2i-y_1^2i+4y_1y_2i-y_2^2i+2x_2y_1-2x_1y_2-8x_2y_2+4z,\\
%    b =& 2x_1^2z^2i+4x_1x_2z^2i+2x_2^2z^2i+2y_1^2z^2i-4y_1y_2z^2i+2y_2^2z^2i-2x_1^2zi-4x_1x_2zi-2x_2^2zi \\ &-2y_1^2zi+4y_1y_2zi-2y_2^2zi+4z^3i-x_1^2i-4x_1x_2i-x_2^2i-y_1^2i+4y_1y_2i-y_2^2i-6z^2i+i. 
%\end{align*}
%
%Over $\real$, there are no ring isomorphism between $k[\mathbb{S}^{n}_\real]$ and $k[Q_{{2n}}]$ or $k[Q_{{2n+1}}]$, but their real points are homotopy equivalent. Wood \william{ref} proved that all maps $\mathbb{S}_\real^4 \to \mathbb{S}_\real^3$ are constant.  
%
%The real realization of the second motivic hopf map $\nu\colon Q_7 \to Q_4$ is $\eta$. Taking the real realization of the $\mathbb{P}^1$-suspension of $\nu$ produces an algebraic representative of the suspension of the hopf map over the reals. 
%
%There's an embedding $\mathbb{S}^{n}_\real \to Q_{2n}$ which is a homotopy equivalence. With this, we can produce a map $\mathbb{S}^{4}_\real \to Q_6$ 
%
%\william{This is all kinda obvious by playing with $\A^1$-naivity, but maybe it's nice to say anyways.}


\subsection*{Conventions}
We work over a perfect field $k$. We use the conventions of \cite{Asok20b} for motivic homotopy theory, and the notation of \cite[\S 3]{Asok12c} for higher Grothendieck-Witt sheaves (a.k.a. Hermitian $K$-theory sheaves). In particular, we will use symplectic $K$-theory and denote the usual skew-symmetric matrix on $\Z^2$ by $H$, i.e.,
\[
H := \begin{pmatrix}
    0 & 1 \\ -1 & 0
\end{pmatrix}.
\]

\section{Maps between quadrics}

\subsection{Basics on quadrics}\label{sec:basicsquadrics}

We start by exploring a bit the geometry of $Q_{2n}$ for $n\geq 1$. We consider the closed subschemes 
\[
Z_0:=\{x_1=\ldots=x_n=z=0\}\subset Q_{2n}, \phantom{i}Z_1:=\{x_1=\ldots=x_n=1-z=0\}\subset Q_{2n}
\]
and set $X_i:=Q_{2n}\smallsetminus Z_i$ for $i=0,1$. It is straightforward to check that $Z_i\simeq \A^n$ (with coordinates $y_1,\ldots,y_n$) and we know from \cite[Theorem 3.1.1]{Asok14} that the subschemes $X_i$ are $\A^1$-contractible. On the other hand
\[
X_0\cap X_1=Q_{2n}\smallsetminus (Z_0\sqcup Z_1)=Q_{2n}\smallsetminus \{x_1=\ldots=x_n=0\}
\]
and there is an obvious projection morphism 
\[
p_n\colon X_0\cap X_1\to \A^n\smallsetminus \{0\},\phantom{i}(x_1,\ldots,x_n,y_1,\ldots,y_n,z)\mapsto (x_1,\ldots,x_n).
\]
\begin{lem}
The map $p_n\colon  X_0\cap X_1\to \A^n\smallsetminus \{0\}$ is a weak equivalence. 
\end{lem}

\begin{proof}
We can cover $\A^n\smallsetminus\{0\}$ by the open subschemes $D(x_i)=\{x_i\neq 0\}\subset \A^n\smallsetminus\{0\}$. The preimage of $D(x_i)$ under $p_n$ is an affine space over $D(x_i)$ with coordinates $y_1,\ldots,y_{i-1},y_{i+1},\ldots,y_n,z$. 
\end{proof}

We can construct an inverse (up to homotopy) of that map, using the morphism
\[
\alpha_n\colon Q_{2n-1}\times \A^1\to Q_{2n}
\]
given on coordinates by $(x_1,\ldots,x_n,y_1,\ldots,y_n,t)\mapsto (x_1,\ldots,x_n,y_1t(1-t),\ldots,y_nt(1-t),t)$. A direct computation shows that it factors through $X_0\cap X_1$, and that the composite 
\[
Q_{2n-1}\times \A^1\xrightarrow{\alpha_n} X_0\cap X_1\xrightarrow{p_n} \A^n\smallsetminus \{0\}
\]
For further use, we also note that $\alpha_n$ induces an isomorphism 
\[
\alpha_n\colon Q_{2n-1}\times (\A^1\smallsetminus \{0,1\})\simeq D(z(1-z))\subset Q_{2n}
\]
with inverse given $(x_1,\ldots,x_n,y_1,\ldots,y_n,z)\mapsto (x_1,\ldots,x_n,\frac {y_1}{z(1-z)},\ldots,\frac {y_n}{z(1-z)},z)$. 

\subsection{The universal rank $2$ bundle on $Q_4$}\label{sec:universal2bundle}

Recall from \cite{Panin21} that there is an isomorphism $Q_4\simeq \mathrm{H}\Proj^1$, where we see the right hand side as the open subscheme of the Grassmannian $\mathrm{Gr}(2,4)$ over which the symplectic form $H\perp H$ restricts to a nondegenerate symplectic form on the universal rank $2$ sub-bundle. We consider the Plücker embedding 
\[
\mathrm{Gr}(2,4)\to \mathbb{P}^5
\]
At a point $x\in  \mathrm{Gr}(2,4)$ given by a matrix
\[
0\to k(x)^2\xrightarrow{\begin{pmatrix} a_1 & b_1 \\ a_2 & b_2 \\ a_3 & b_3 \\ a_4 & b_4\end{pmatrix}} k(x)^4
\]
the condition to be non degenerate is given by $a_1b_2-a_2b_1+a_3b_4-a_4b_3\neq 0$. Using the Plücker coordinates 
\[
(u_1,u_2,u_3,u_4,u_5,u_6)=(a_1b_2-a_2b_1,a_1b_3-a_3b_1,a_1b_4-a_4b_1,a_2b_3-a_3b_2,a_2b_4-a_4b_2,a_3b_4-a_4b_3)
\]
this corresponds to $u_1+u_6\neq 0$. Setting $d:=u_1+u_6$, and using the coordinates 
\[
z:=\frac {a_1b_2-a_2b_1}d,x_1:=\frac {a_1b_3-a_3b_1}d, x_2:=\frac {a_1b_4-a_4b_1}d, y_1:=\frac {a_2b_4-a_4b_2}d, y_2:=\frac {a_3b_2-a_2b_3}d
\]
we obtain that the (unique) equation describing $\mathrm{Gr}(2,4)\cap D_+(d)$ is given by $x_1y_1+x_2y_2=z(1-z)$. 

Using this, we can describe the canonical rank $2$ bundle (and its canonical rank $2$ orthogonal) on $Q_4\simeq \mathrm{H}\Proj^1$. We will do it under the form of an idempotent matrix, that we now make explicit. We consider the commutative diagram
\[
\xymatrix{P\ar[r]^-{i}\ar[d]_-{\varphi} & R^4\ar[d]^-{H\perp H} \\
P^\vee & (R^4)^\vee\ar[l]^-{i^\vee}}
\]
where $R$ is the coordinate ring of $Q_4$, $P$ is the universal rank $2$ bundle and $\varphi\colon P\to P^\vee$ is the symplectic form induced by $H\perp H$. We then get an idempotent matrix using $i\varphi^{-1}i^\vee(H\perp H)$. 

At a point $x\in  \mathrm{Gr}(2,4)$ given by the above matrix $\begin{pmatrix} a_1 & b_1 \\ a_2 & b_2 \\ a_3 & b_3 \\ a_4 & b_4\end{pmatrix}$, a direct computation yields a matrix of the form
\[
\frac 1{d}\begin{pmatrix} a_1b_2-a_2b_1 & 0 & a_1b_4-a_4b_1 & a_3b_1-a_1b_3 \\ 0 & a_1b_2-a_2b_1 & a_2b_4-a_4b_2 & a_3b_2-a_2b_3 \\
a_3b_2-a_2b_3 & a_1b_3-a_3b_1 & a_3b_4-a_4b_3 & 0 \\ a_4b_2-a_2b_4 & a_1b_4-a_4b_1 & 0 & a_3b_4-a_4b_3\end{pmatrix}
\]
which corresponds to the matrix
\[
M:=\begin{pmatrix} z & 0 & x_2 & -x_1 \\ 0 & z & y_1 & y_2 \\
y_2 & x_1 & 1-z & 0 \\ -y_1 & x_2 & 0 & 1-z\end{pmatrix}.
\]
We let the reader check that this module coincides with the one in \cite[Example 4.3.8]{Asok14}.
The orthogonal of the rank $2$ universal subbundle is then given by the matrix
\[
N:=\mathrm{Id}-M=\begin{pmatrix} 1-z & 0 & -x_2 & x_1 \\ 0 & 1-z & -y_1 & -y_2 \\
-y_2 & -x_1 & z & 0 \\ y_1 & -x_2 & 0 & z\end{pmatrix}.
\]

\subsection{The Jouanolou device of $\mathbb{P}^n$}

Recall that the usual Jouanolou device $\mathbb{J}^n$ of $\Proj^n$ is given by the functor associating to a $k$-algebra $R$ the set of matrices $P\in M_{n+1}(R)$ of rank one satisfying $P^2=P$. We have an explicit presentation of the algebra representing this functor via
\[
k[\mathbb{J}^n]=k[x_{ij}\vert 1\leq i,j\leq n+1]/\left ( x_{ij}=\sum_{r=1}^{n+1} x_{ir}x_{rj}\forall 1\leq i,j\leq n+1, D\right)
\]
where $D$ is the set of all $2\times 2$ minors of $P$.
There is a universal matrix $P=(x_{ij})$ on $k[\mathbb{J}^n]$ which satisfies $P^2=P$ by the first set of equations. This yields a projector $P_R$ on $R^{n+1}$ for any homomorphism of $k$-algebras $k[\mathbb{J}^n]\to R$ and we set $L=\mathrm{Im}(P)\subset R^{n+1}$. This is a projective module and we can compute its rank by considering residue fields of $R$. If $F$ is a field we know that the matrix $P_F$ satisfies $P^2_F=P_F$ and it follows in particular that $P_F$ is diagonalizable with possible eigenvalues $0$ and $1$. The condition on the minors forces $P_F$ to be of rank one, and consequently $L$ is of rank one as required. 

\begin{rem}
If $k$ is of characteristic $0$, we can replace the conditions $D=0$ by the easier condition $\sum_{i=1}^{n+1}x_{ii}=1$, i.e. the matrix is of trace $1$.  
\end{rem}

We have an obvious morphism of schemes 
\[
p:\mathbb{J}^n\to \Proj^n
\]
corresponding to the line bundle $L$ and its generating sections given by $k[\mathbb{J}^n]^{n+1}\xrightarrow{P} L$. Concretely, let $c_1,\ldots,c_{n+1}$ (resp. $l_1,\ldots,l_{n+1}$) be the columns (resp. rows) of $P$. We set $U_i=\Jo^n\smallsetminus V(c_i)$, where $V(c_i):=V(x_{1i},\ldots,x_{(n+1)i})$ and we observe that $c_i$ is non-trivial on $U_i$, i.e. $c_i$ generates $L_{\vert U_i}$. On $U_i$, there exists then global sections $\lambda_{i,j}$ for $j=1,\ldots,n+1$ with $c_j=\lambda_{i,j}c_i$ and we obtain a map
\[
p_i:U_i\to \A^n
\]
given by the global sections $\lambda_{i,j}$, $i\neq j$. These maps glue together to yield the morphism $p:\mathbb{J}^n\to \Proj^n$. It turns out that $U_i$ is in fact isomorphic to $\A^{2n}$. To see this, we observe that we can write 
\[
P_{\vert U_i}=c_i\cdot \begin{pmatrix} \lambda_{i,1} & \ldots & \lambda_{i,i-1} & 1 & \lambda_{i,i+1} & \ldots & \lambda_{i,n+1}\end{pmatrix}.
\]
From $P^2=P$, we draw in particular that $P\cdot c_i=c_i$ and thus 
\[
\begin{pmatrix} \lambda_{i,1} & \ldots & \lambda_{i,i-1} & 1 & \lambda_{i,i+1} & \ldots & \lambda_{i,n+1}\end{pmatrix}\cdot c_i=1
\]
which reads as 
\[
\sum_{j=1, j\neq i}^{n+1}\lambda_{i,j}x_{ji}+x_{ii}=1.
\]
and $x_{ii}$ is determined by the other elements. This allows to define an isomorphism $\varphi_i:\A^{2n}\to U_i$ via
\[
(a_1,\ldots,a_n,b_1,\ldots,b_n)\mapsto \begin{pmatrix} a_1 \\ \vdots \\ a_{i-1} \\ 1-\sum_{j=1}^na_jb_j \\ a_i \\ \vdots \\ a_n\end{pmatrix}\cdot 
\begin{pmatrix} b_1 & \cdots & b_{i-1} & 1 & b_i & \cdots & b_n\end{pmatrix}.
\]
The point $0:=[0:\ldots:0:1]\in \mathbb{P}^n$ is contained in the image of the projection 
\[
p_{n+1}\colon U_{n+1}\to \A^n\subset \mathbb{P}^n
\]
and using $\varphi_{2n+1}$ we see that its preimage is of the form
\[
p_{n+1}^{-1}(\{0\})=\{ \begin{pmatrix} 0 & \hdots & 0 & a_1 \\ \vdots & \ddots & \vdots & \vdots \\ 0 & \hdots & 0 & a_n \\ 0 & \hdots & 0 & 1\end{pmatrix}\vert (a_1,\ldots,a_n)\in \A^n\}\simeq \A^n.
\] 
and we obtain a Cartesian square 
\begin{equation}\label{eqn:compatiblestructures}
\xymatrix{\A^n\ar[r]\ar[d] & \mathbb{J}^n\ar[d] \\
\{0\}\ar[r] & \mathbb{P}^n}
\end{equation}



\subsection{The map $\mathbb{J}^n\to Q_{2n}$}\label{sec:collapsemap}

Consider the ideal 
\[
I=\langle x_{ij}\vert 1\leq i\leq n+1, 1\leq j\leq n\rangle\subset k[\mathbb{J}^n].
\]
We claim that 
\[
I=\langle x_{(n+1)1},\ldots x_{(n+1),n},\sum_{i=1}^n x_{ii}\rangle=\langle x_{(n+1)1},\ldots x_{(n+1),n},1-x_{(n+1),(n+1)}\rangle=p_{n+1}^{-1}(\{0\}).
\]
 Indeed, for $1\leq i,j\leq n$ the $2\times 2$-minor
\[
\det \begin{pmatrix} x_{ij} & x_{i(n+1)} \\ x_{(n+1)j} & x_{(n+1)(n+1)}\end{pmatrix}
\]
is trivial, yielding
\[
x_{ij}(1-\sum_{i=1}^nx_{ii})=x_{ij}x_{(n+1)(n+1)}=x_{i(n+1)}x_{(n+1)j}.
\]
Thus $x_{ij}=x_{i(n+1)}x_{(n+1)j}+x_{ij}(\sum_{i=1}^nx_{ii})$ and the claim follows. Using again
\[
\det \begin{pmatrix} x_{ii} & x_{i(n+1)} \\ x_{(n+1)i} & x_{(n+1)(n+1)}\end{pmatrix}=0
\]
and summing up, we find
\[
\sum_{i=1}^nx_{i(n+1)}x_{(n+1)i}=(\sum_{i=1}^nx_{ii})x_{(n+1)(n+1)}=(\sum_{i=1}^nx_{ii})(1-\sum_{i=1}^nx_{ii})
\]
and it follows that $I/I^2$ is generated by the classes of $x_{(n+1)1},\ldots x_{(n+1),n}$. This yields a morphism 
\[
\pi_n\colon \mathbb{J}^n\to Q_{2n}
\] 
given by $(x_{ij})\mapsto (x_{(n+1)1},\ldots x_{(n+1),n},x_{1(n+1)},\ldots,x_{n(n+1)},\sum_{i=1}^nx_{ii})$. Moreover, the Cartesian square \eqref{eqn:compatiblestructures} yields a Cartesian square
\begin{equation}\label{eqn:commutativeinvariant}
\xymatrix{\mathbb{J}^n\smallsetminus \{\A^n\}\ar[r]\ar[d] & \mathbb{J}^n\ar[d] \\ 
\mathbb{P}^n\smallsetminus \{0\} \ar[r] & \mathbb{P}^n}
\end{equation}
the induced morphism on cofibers (using the purity isomorphism) reads as a morphism 
\[
\mathbb{J}^n\to (\A^n)_+\wedge (\pone)^{\wedge n}\to (\pone)^{\wedge n}
\]
for which $\pi_n$ is an explicit model.

\subsection{An explicit $\pone$-suspension map of maps between quadrics}
\label{sec:pone-suspension}
In this section, we provide an explicit way of $\pone$-suspending morphisms of quadrics, starting with the following observation: For $n\geq 1$, consider  $\A^{2n+2}$ with variables $x_1,\ldots,x_n,y_1,\ldots,y_n,t,u$. We have $f_n:=\sum_{i=1}^n x_iy_i-1$ and $g=tu-1$. These polynomials are irreducible, hence prime in $k[\A^{2n+2}]$, and
\[
V(f_n)=Q_{2n-1}\times \A^2, \phantom{i}V(g)=\A^{2n}\times \gm{}.
\]  
Since $f_n$ and $g$ are prime, we have $(f_n)\cap (g)=(f_n\cdot g)$ and obtain in view of \cite[Thm. 3.4, Thm. 3.5, Cor. 3.9]{Schwede05} a cocartesian square of schemes (whose arrows are the closed immersions given by the above polynomials)
\[
\xymatrix{Q_{2n-1}\times \gm{}\ar[r]\ar[d] & Q_{2n-1}\times \A^2\ar[d] \\
\A^{2n}\times \gm{}\ar[r] & X_{2n+1}}
\]
where $X_{2n+1}=V(f_n\cdot g)\subset \A^{2n+2}$. Arguing as in \cite[Section 3, Lemma 1.6]{Morel99}, it is not hard to see that $X_{2n+1}$ is, in fact, a push-out in the category of Nisnevich sheaves of sets on $\mathrm{Sm}_k$. In fact, we may use \cite[\S 7.3]{Morel08} to see that it is a model of the reduced join of the corners of the diagram, which is  $\Sigma Q_{2n-1}\wedge \gm{}\simeq Q_{2n+1}$. It is slightly inconvenient to work with singular schemes, and we now provide an explicit weak equivalence $X_{2n+1}\to  Q_{2n+1}$. We have morphisms 
\[
\varphi_1:Q_{2n-1}\times \A^2\to Q_{2n+1}
\]
given by $(x_1,\ldots,x_n,y_1,\ldots,y_n,t,u)\mapsto (x_1,\ldots,x_n,t,y_1,\ldots,y_n,0)$
and 
\[
\varphi_2:\A^{2n}\times \gm{}\to Q_{2n+1}
\]
defined by $(x_1,\ldots,x_n,y_1,\ldots,y_n,t,u)\mapsto (x_1,\ldots,x_n,t,y_1,\ldots,y_n,u(1-\sum_{i=1}^nx_iy_i))$. It is straightforward to check that these maps induce a map 
\[
\varphi:X_{2n+1}\to Q_{2n+1}.
\]
Concretely, the morphism $\varphi$ is given by
\[
k[x_1,\ldots,x_{n+1},y_1,\ldots,y_{n+1}]/f_{n+1} \to k[x_1,\ldots,x_n,y_1,\ldots,y_n,u,t]/(f_n\cdot g)
\]
with $(x_1,\ldots,x_n,y_1,\ldots,y_n)\mapsto (x_1,\ldots,x_n,y_1,\ldots,y_n)$, $x_{n+1}\mapsto t$ and $y_{n+1}\mapsto uf_n$.  To prove that it is a weak equivalence, we observe that the projection $\pi\colon Q_{2n+1}\to \A^{n+1}\smallsetminus\{0\}$ onto the first $n+1$-coordinates is a weak equivalence since its fibers are affine spaces of dimension $n$.  We can consider the preimages of the open subschemes $\A^n\times \gm{}$ and $(\A^n\smallsetminus\{0\})\times \A^1$ denoted respectively by $D(x_{n+1})$ and $U_n$. For the same reasons as above, the projections $\pi\colon D(x_{n+1})\to \A^n\times \gm{}$ and $\pi\colon U_n\to (\A^n\smallsetminus\{0\})\times \A^1$ are weak equivalences, and the same holds on the intersection $U_n\cap D(x_n)$ which becomes weakly equivalent to $(\A^n\smallsetminus\{0\})\times \gm{}$. Now, the morphism 
\[
\varphi_1\colon Q_{2n-1}\times \A^2\to Q_{2n+1}
\]
factorizes through $U_n$ by definition. This induces a weak equivalence $\varphi_1\colon  Q_{2n-1}\times \A^2\to U_n$ since its composite with the projection  $\pi\colon U_n\to (\A^n\smallsetminus\{0\})\times \A^1$ corresponds to $(x_1,\ldots,x_n,y_1,\ldots,y_n,t,u)\mapsto (x_1,\ldots,x_n,t)$. The same argument shows that the morphism
\[
\varphi_2:\A^{2n}\times \gm{}\to Q_{2n+1}
\]
also induces a weak equivalence $\A^{2n}\times \gm{}\to D(x_{n+1})$, while both morphisms induce weak equivalences $ Q_{2n-1}\times\gm{}\to  U_n\cap D(x_{n+1})$. Using  \cite[Section 3, Lemma 1.6]{Morel99}, we obtain a push-out square in $\mathcal{H}(k)$ of the form
\[
\xymatrix{U_n\cap D(x_{n+1})\ar[r]\ar[d] &  D(x_{n+1})\ar[d] \\
 U_n\ar[r] & Q_{2n+1}}
\] 
and the previous arguments show that $\varphi\colon X_{2n+1}\to Q_{2n+1}$ is a weak equivalence.


We proceed in an analogous manner for $Q_{2n}$. The details are as follows. For $n\geq 1$, consider $\A^{2n+3}$ with variables $x_1,\ldots,x_n,y_1,\ldots,y_n,z,t,u$. Set $h_n=\sum_{i=1}^nx_iy_i-z(1-z)$ and $g=tu-1$. The same argument as before yields a cocartesian square 
\[
\xymatrix{Q_{2n}\times \gm{}\ar[r]\ar[d] & Q_{2n}\times \A^2\ar[d] \\
\A^{2n+1}\times \gm{}\ar[r] & Y_{2n+2}}
\]
where $Y_{2n+2}=V(h_n\cdot g)\subset \A^{2n+3}$. We have morphisms
\[
\psi_1:Q_{2n}\times \A^2\to Q_{2n+2}
\]
given by $(x_1,\ldots,x_n,y_1,\ldots,y_n,z,t,u)\mapsto (x_1,\ldots,x_n,t,y_1,\ldots,y_n,0,z)$
and 
\[
\psi_2:\A^{2n+1}\times \gm{}\to Q_{2n+2}
\]
defined by $(x_1,\ldots,x_n,y_1,\ldots,y_n,z,t,u)\mapsto (x_1,\ldots,x_n,t,y_1,\ldots,y_n,u(z(1-z)-\sum_{i=1}^nx_iy_i),z)$. They induce a morphism 
\[
\psi:Y_{2n+2}\to Q_{2n+2}
\]
which is concretely given by 
\[
k[x_1,\ldots,x_{n+1},y_1,\ldots,y_{n+1},z]/h_{n+1}\to k[x_1,\ldots,x_{n},y_1,\ldots,y_{n},z,t,u]/(h_{n}\cdot g)
\]
with $(x_1,\ldots,x_{n},y_1,\ldots,y_{n},z)\mapsto (x_1,\ldots,x_{n},y_1,\ldots,y_{n},z)$, $x_{n+1}\mapsto t$ and $y_{n+1}\mapsto u(z(1-z)-\sum_{i=1}^nx_iy_i)$.

To see that $\psi$ is a weak equivalence, observe that we have open contractible subschemes $X_i\subset Q_{2n}$ for $i=0,1$. The morphism
\[
\psi_1:Q_{2n}\times \A^2\to Q_{2n+2}
\]
induces a morphism $X_0\times \A^2\to Q_{2n+2}$, which is easily seen to factor through the contractible subscheme $X_0':=D(x_1,\ldots,x_{n+1},z)\subset Q_{2n+2}$. Similarly, the morphism 
\[
\psi_2:\A^{2n+1}\times \gm{}\to Q_{2n+2}
\]
factors through $X_0'$ and we obtain a morphism from the push-out $P_0$ of the diagram
\[
\xymatrix{X_0\times \gm{}\ar[r]\ar[d] & X_0\times \A^2 \\
\A^{2n+1}\times \gm{} &  }
\]
to $X_0'$. Since $X_0$ and $\A^{2n+1}$ are both contractible, it is straightforward to check that the push-out is contractible and thus the induced map $Q_0\to X_0'$ is a weak equivalence. The same arguments apply replacing $X_0$ by $X_1$ obtaining a push-out $P_1$. Taking the intersection we thus obtain a morphism of the push-out $P_{01}$ of the diagram 
\[
\xymatrix{(X_0\cap X_1)\times \gm{}\ar[r]\ar[d] & (X_0\cap X_1)\times \A^2 \\
\A^{2n+1}\times \gm{} & }
\]
to $X'_{0}\cap X'_1$. The push-out of the diagram involving $P_{01},P_0,P_1$ is $Y_{2n+2}$, while the push-out of the diagram involving $X'_0\cap X'_1,X'_0,X'_1$ is $Q_{2n+2}$. This proves that the map from $Y_{2n+2}$ to $Q_{2n+2}$ is a weak equivalence. Once again, we think about $Y_{2n+2}$ as the singular $\pone$-suspension of $Q_{2n}$.

The construction described above allows to explicitly construct the $\pone$-suspension of a morphism $f\colon Q_{2n-1}\to Q_{2m}$ defined by mapping $(x_1,\ldots,x_n,y_1,\ldots,y_n)$ to $(f_1, \ldots, f_n, \tilde{f_1},\ldots, \tilde{f}_n, f_z)$. First note that an equivalent condition for $f$ being a morphism $Q_{2n-1}\to Q_{2m}$ is that the ideal 
\[
\langle f_1, \ldots, f_n, \sum_{i=1}^n x_iy_i -1 \rangle
\] 
in the ring $k[\A^{2n}]$ contains the element $f_z(1-f_z)$.

One observes that $f$ lifts to a map $F:\A^{2n}\to \A^{2m+1}$ given by the same formula, and to a map $(F)\times \mathrm{Id}_2:\A^{2n+2}\to \A^{2m+3}$ which descends to a map
\[
X_{2n+1}\to Y_{2m+2}
\]
and thus to a map $X_{2n+1}\to Q_{2m+2}$ which is given by mapping $(x_1,\ldots,x_n,y_1,\ldots,y_n,t,u)$ to 
\[
(f_1,\ldots f_n,t, \tilde f_1 ,\ldots\tilde f_n,\alpha,f_z)
\]
where $\alpha$ is an element in $k[X_{2n}]$ such that $\sum f_i\tilde f_i + t\alpha = f_z(1-f_z)$. Since $f\colon Q_{2n-1} \to Q_{2m}$ is a morphism, there exists an element $r$ such that $\sum f_i \tilde f_i + r(\sum x_i y_i -1) = f_z(1-f_z)$ in $k[\A^{2n}]$. By letting $\alpha = ur(\sum x_i y_i -1)$, we see that 
$\sum f_i\tilde f_i + t\alpha - r(tu-1)(\sum x_i y_i -1) = f_z(1-f_z)$ in $k[\A^{2n+2}].$

We then are left to find a factorization 
\[
X_{2n+1}\to Q_{2n+1}\to Q_{2m+2}
\]
of the above map. With a bit of thought, we realize that we may consider 
\[
Q_{2n+1}\to Q_{2m+2}
\]
given by the mapping $(x_1,\ldots, x_{n+1}, y_1, \ldots, y_{n+1})$ to 
\[
(f_1, \ldots, f_n, x_{n+1},\tilde{f}_1, \ldots,\tilde f_n, y_{n+1}r,  f_z)
\]
which does the job.

Doing the same arguments for maps $Q_{2n-1}\to Q_{2m-1}, Q_{2n}\to Q_{2m-1}$ and $Q_{2n}\to Q_{2m}$ give us the following results.

\begin{lem}
\label{lem:ponesuspoddeven}
The $\pone$-suspension of a map $f:Q_{2n-1}\to Q_{2m}$ (resp.$f:Q_{2n}\to Q_{2m})$  given by the ideal $\langle f_1, \ldots, f_n, f_z \rangle $ is the map $F:Q_{2n+1}\to Q_{2m+2}$ (resp. $F:Q_{2n+2}\to Q_{2m+2}$) given by the ideal $\langle f_1, \ldots, f_n, x_{n+1},f_z \rangle$.
\end{lem}


\begin{lem}
\label{lem:ponesuspevenodd}
The $\pone$-suspension of a map $f:Q_{2n}\to Q_{2m-1}$ (resp. $f:Q_{2n-1}\to Q_{2m-1}$) given by the unimodular row $(f_1, \ldots, f_n)$ is the map $F:Q_{2n+2}\to Q_{2m+1}$ (resp. $F:Q_{2n+1}\to Q_{2m+1}$) given by the unimodular row $(f_1, \ldots, f_n, x_{n+1})$.
\end{lem}

\begin{rem}
The maps of the two above lemmas are actually pointed, if we choose $(1,0\ldots,0,1,0,\ldots,0)$ as the base point of $Q_{2n+1}$ and $(0,\ldots,0)$ as the base point of $Q_{2m}$ in the first case, and $(1,0\ldots,0,1,0,\ldots,0)$ as the base point of $Q_{2n+1}$ and $(0,\ldots,0)$ as the base point of $Q_{2m}$ in the second case.
\end{rem}
\subsection{Endomorphisms of quadrics of arbitrary $\A^1$-Brouwer degree}
In this section we give a procedure for constructing morphisms $Q_n \to Q_n$ representing each element in $\mathrm{GW}(k)$, for $n\geq 2$. We begin by constructing the morphisms in the cases $[Q_2,Q_2]_{\A^1}$ and $[Q_3,Q_3]_{\A^1}$ and then suspend accordingly using Lemma \ref{lem:ponesuspevenodd} and \ref{lem:ponesuspoddeven}. We start with considering endomorphisms of $Q_3$. 
\begin{lem}
    The map $Q_3 \to Q_3$ given by the unimodular row $(x_1, ux_2)$ represents the class $\langle u \rangle \in \mathrm{GW}(k)$ and the row $(y_1, -ux_2)$ represents the class $-\langle u \rangle \in \mathrm{GW}(k)$
\end{lem}

\begin{proof}
    For a unit $u\in k^\times$, the map $\A^2 \smallsetminus \{0\} \to \A^2 \smallsetminus \{0\}$ given on coordinates by $(x_1, x_2) \mapsto (x_1,ux_2)$ has $\A^1$-Brouwer degree $\langle u \rangle $ \cite[Lemma 9]{Kass19}. Since $Q_3 \simeq \A^2 \smallsetminus \{0\}$, this endomorphism lifts to the map $Q_3 \to Q_3$ given by the $\mathrm{SL}_2(k[Q_3])$-matrix. \[
    m_u = \begin{pmatrix}
        x_1 & ux_2\\ -u^{-1}y_2  & y_1
    \end{pmatrix}.
    \]
    The group structure on $[Q_3,Q_3]_{\A^1}$ is given by the group structure on $\mathrm{SL}_2$, hence the morphism representing $-\langle u\rangle $ is given by the inverse matrix 
    \[
   m_u^{-1} =  \begin{pmatrix}
        x_1 & ux_2\\ -u^{-1}y_2  & y_1
    \end{pmatrix}^{-1} =\begin{pmatrix}
        y_1 & -ux_2\\ u^{-1}y_2  & x_1
    \end{pmatrix} .
    \]
\end{proof}
Given an element in $q \in \mathrm{GW}(k)$, one does the following to find a morphism $Q_3 \to Q_3 $ with $\A^1$-Brouwer degree $q$:
\begin{enumerate}
    \item Write $q$ as a sum $q = \langle u_1, \ldots, u_r\rangle - \langle v_1, \ldots, v_s\rangle$. 
    \item Compute the matrix product $M_q = m_{u_1} \ldots m_{u_r} m_{v_1}^{-1} \ldots m_{v_s}^{-1}.$
    \item The first row of $M_q$ is a unimodular row with $\A^1$-Brouwer degree $q.$
\end{enumerate}
We will now consider $Q_2$-endomorphisms. Note that $Q_2 \simeq \mathbb{J}^1\simeq \pone$ and that in \cite{Barth25} the group structure on $[Q_2,\Proj^1]_{\A^1}$ is studied. The first complication is that $[Q_2,Q_2]_{\A^1}$ is not in the stable range, and we have $[Q_2,Q_2]_{\A^1} \cong \mathrm{GW}(k) \times _{k^\times / k^{\times 2}} k^\times$, but under stabilization, this map is just projection onto $\mathrm{GW}(k)$ (use for instance \cite[Theorem 6.37]{Morel08}. The group operation on $[Q_2,\Proj^1]_{\A^1}$ is in general not easy to describe, but the operation of adding a rank $0$ map to any other map can be described easily. We first begin by describing the generators of the rank $0$ maps. 
\begin{lem} 
\label{lem:q2rank0maps}
    Let $u,v \in k^\times$, then the stable class of the morphism $f_{u,v}\colon Q_2 \to Q_2 $ given by the ideal $\langle \frac{(u-v)}{uv}x_1,  z + \frac{v}{u}(1-z)\rangle$ is $\langle u\rangle - \langle v\rangle$. 
\end{lem}
\begin{proof}
     It follows from \cite[Proposition 103]{Barth25} that the stable class of the ideal $\langle ux_1,  z\rangle$ is $\langle u\rangle$, and from \cite[Lemma 91]{Barth25} that the class of $f_{u,v} \oplus f_v$ is $\langle u\rangle$. Hence the class of $f_{u,v}$ is $\langle u\rangle - \langle v\rangle $. 
\end{proof}
Since all $Q_2$-endomorphisms representing rank $0$ forms factor through $Q_3$ (\cite[Corollary 43]{Barth25}), the ideal described in Lemma \ref{lem:q2rank0maps} is in fact a unimodular row. Thus we define the $\mathrm{SL_2}(k[Q_2])$-matrix \[m_{u,v} := \begin{pmatrix}
    z + \frac{v}{u}(1-z) &  \frac{(u-v)}{uv}x_1 \\ (u-v)y_1  & z + \frac{u}{v}(1+z) 
\end{pmatrix}.\]
Consider polynomials $A_n(z), B_n(z)$ defined such that $(z + (1-z))^{2n} = z^n A_n(z) + (1-z)^n B_n(z)$. 
\begin{lem}
\label{lem:Q2hyperbolic}
    Let $n>0$ then: 
    \begin{enumerate}
        \item The map $Q_2 \to Q_2$ given by the ideal $\langle z^n A(n), x_1^nA_n(z)\rangle$ has degree 
    $\frac{n}{2}\langle 1,-1\rangle$ when $n$ is even and  $\frac{n-1}{2}\langle 1,-1\rangle + \langle1\rangle$ when $n$ is odd. 
    
    \item The map $Q_2 \to Q_2$ given by the ideal $\langle z^n A_n(z), y_1^nA_n(z)\rangle$ has degree 
    $\frac{n}{2}\langle 1,-1\rangle$ when $n$ is even and  $-\frac{n-1}{2}\langle 1,-1\rangle - \langle -1\rangle$ when $n$ is odd.
    \end{enumerate}
\end{lem}
\begin{proof}
    Since these morphisms are defined over $\Z$, one can use a similar argument as the proof of Theorem 105 in \cite{Barth25} to show that the class the morphisms are determined by the Brouwer degree of their real and complex realizations. 
\end{proof}
 A $Q_2$-endomorphism can be described as a $(2 \times 2)$ idempotent matrix $P$ with trace $1$ and entries in $k[Q_2]$. For $n \geq 0$, we define the idempotent matrices
 \[
 P_n = \begin{pmatrix}
     z^n A_n(z) & x_1^nA_n(z) \\ y_1^n B_n(z) & (1-z)^nB_n(z)
 \end{pmatrix} \quad P_{-n}= \begin{pmatrix}
     z^n A_n(z) & y_1^nA_n(z) \\ x_1^n B_n(z) & (1-z)^nB_n(z)
 \end{pmatrix}.\]
In loc.cit. an action of rank $0$ maps on rank $n$ maps is described, and it is shown that this agrees with the group operation on $[Q_2, \pone]^{\A^1}$. One can check, following Remark 84 loc.cit, that the action of $m_{u,v}$ on $P$ is given by the matrix $m_{u,v}Pm_{u,v}^{-1}$.

Given an element in $q \in \mathrm{GW}(k)$, one does the following to find a morphism $Q_2 \to Q_2 $ that stably represents $q$. 
\begin{enumerate}
    \item If $q$ is of even rank $n$, write it as a sum $\frac{n}{2}\langle 1, -1\rangle + \sum_{i=1}^r \langle u_i\rangle - \langle v_i\rangle$ for some units $u_i, v_i \in k^\times$. If $q$ is of odd rank $n$, write it as a sum $\frac{n-1}{2}\langle 1, -1\rangle +\langle 1\rangle  + \sum_{i=1}^r \langle u_i\rangle - \langle v_i\rangle$ for some units $u_i, v_i \in k^\times$.
    \item One then computes matrix product $P_q= m_{u_1,v_1}\ldots m_{u_r,v_r}P_n m_{u_r,v_r}^{-1}\ldots m_{u_1,v_1}^{-1}$
    \item The first column of $P_q$ is the ideal defining the map $Q_2 \to Q_2$ with stable class $q$. 
\end{enumerate}
\begin{rem}
    The reason we do not decompose $q$ as a sum $n\langle 1\rangle + \sum_{i=1}^m \langle u_i\rangle - \langle v_i\rangle$ is that the $Q_2$-endomorphism representing $n\langle 1\rangle $ is a lot more complicated to write down compared to the maps given in Lemma \ref{lem:Q2hyperbolic}.
\end{rem}
\subsection{Motivic Hopf maps}

Given an H-space $X$, the Hopf construction detailed in \cite[\S 7.3]{Morel08} provides us with a fiber sequence 
\begin{equation*}
    X \to X\star X \to \Sigma X. 
\end{equation*}
Applying this to the algebraic group $\gm$ yields the fiber sequence
\begin{equation*}
    \gm{} \to Q_3 \to Q_2.
\end{equation*}
The map $\eta \colon Q_3 \to Q_2$ can be described as the rank $1$ projective $k[Q_3]$-module 
\begin{equation*}
    \begin{pmatrix}
        x_1y_1 & x_2y_1 \\ x_1y_2 & x_2y_2
    \end{pmatrix}.
\end{equation*}

\begin{prop}
\label{prop:etasusp}
Let $n\geq 1$. The map $f\colon Q_{2n+1} \to Q_{2n}$ defined by the ideal $(x_3, \ldots, x_n, x_1x_2)$ is the $(n-1)$-fold $\mathbb{P}^1$-suspension of the Hopf map $\eta \colon Q_3 \to Q_2.$


\end{prop}
\begin{proof}
We first need to prove that the claimed map is the Hopf map when $n=1$. Recall that we have a homotopy equivalence $\pi \colon Q_3 \to \A^2 \smallsetminus\{0\}$ sending $(x_1,x_2,y_1,y_2)$ to the point $(x_1,x_2)$ in $\A^2 \smallsetminus\{0\}$. There is also a homotopy equivalence $Q_2 \to \pone$ which maps a point $(x_1,y_1,z) \in Q_2$ to the point $[z:x_1]$ or $[y_1 : 1-z]$ (whichever is nonzero). The Hopf map $\eta\colon \A^2 \smallsetminus\{0\} \to \pone$ is defined by $(x_1,x_2) \to [x_1:x_2]$. It is straightforward to check that the square 
\begin{equation*}
    \xymatrix{Q_3\ar[r]^\eta\ar[d]^\pi & Q_2\ar[d]^\pi \\
\A^2 \smallsetminus\{0\} \ar[r]_{\eta} & \mathbb{P}^1.}
\end{equation*}
commutes. The claim then follows from Lemma \ref{lem:ponesuspoddeven}.

\end{proof}
Using the Hopf construction for the algebraic group $\mathrm{SL}_2$ yields the fiber sequence
\begin{equation*}
    \mathrm{SL_2} \to Q_7 \xrightarrow{\nu} Q_4. 
\end{equation*}
As explained in \cite[\S3]{Fasel16a}, the scheme $Q_7$ can be seen of the scheme of pairs of $(2\times 2)$-matrices $(M_1,M_2)$ with $\det M_1 + \det M_2 = 1$. The map $\nu\colon Q_7 \to Q_4$ is given by $(M_1 M_2, \det(M_1))$, and its real realization is actually the Hopf map by \cite[Proposition 3.0.1]{Fasel16a}. Note that the complex realization of $\nu$ is the second Hopf map by \cite{Gluck86}. Using once again Lemma \ref{lem:ponesuspoddeven}, we obtain the following proposition.

\begin{prop}\label{prop:nususp}
Let $n\geq 2$. The map $f\colon Q_{2n+3} \to Q_{2n}$ defined by the ideal 
\begin{equation*}
    (-x_3y_1-x_2y_4,-x_4y_1+x_2y_3, x_5, \ldots, x_{n+2}, x_3y_3 + x_4y_4)
\end{equation*} 
is the $(n-2)$-fold $\mathbb{P}^1$-suspension of the Hopf map $\nu \colon Q_7 \to Q_4$.
\end{prop}

Complex realization of the maps from Proposition \ref{prop:etasusp} gives us complex algebraic representatives of the maps $\Sigma^{2n}\eta$ in the category of topological spaces. Similarly, the real realization of the maps from Proposition \ref{prop:nususp} gives us real algebraic representatives of the maps $\Sigma^n\eta$. Our two families of maps do not allow us to create a map whose complex realization is $\Sigma \eta$. Such a map would need to be a map $Q_4 \to Q_3$, and we construct such a map in the next section. 

\section{Exotic Hopf maps}

\subsection{First construction: Symplectic $K$-theory}\label{sec:symplecticHopfmap}
In this section, we construct a map $Q_4 \to Q_3$ that (complex) realizes to the suspension of the Hopf map.  We first note that we have an exact sequence of sheaves \cite[Theorem 3]{Asok12a} 
\begin{equation}\label{eqn:pi2a2minus0}
\KMW4/12_\epsilon\to \piaone_{2}(Q_3)\to \mathbf{GW}_3^2\to 0,
\end{equation}
where $\mathbf{GW}_3^2$ is the Nisnevich sheafification of the presheaf $X\mapsto \mathrm{GW}_3^2(X)=\mathrm{KSp}_3(X)$. Contracting twice and using \cite[Proposition 2.9]{Asok12b}, \cite[Proposition 4.4]{Asok12a} we obtain an exact sequence
\[
\KMW2/12_\epsilon\to \piaone_{2,2}(Q_3)\to \mathbf{GW}_1^0\to 0.
\]
For any field $F/k$
\[
\mathbf{GW}_1^0(F)=\mathrm{KO}_1(F)=F^\times/(F^{\times})^2\times \Z/2\Z
\]
with generator for the $\Z/2\Z$ factor given by the permutation matrix $\begin{pmatrix}0 & 1 \\ 1 & 0\end{pmatrix}$. 


On the other hand, $\mathbf{GW}_3^2=\piaone_2(\mathrm{Sp})$ and the right-hand map in \eqref{eqn:pi2a2minus0} is induced by the inclusion $Q_3\to \mathrm{Sp}$ (recall that $Q_3 = \mathrm{SL}_2 = \mathrm{Sp}_2$). Any map $Q_4\to Q_3$ can be composed with this inclusion, and conversely using the usual fiber sequences and connectivity arguments (e.g. \cite[Corollary 2.4]{Asok12a}), we see that any map $Q_4\to Q_3$ comes from a map $Q_4\to \mathrm{Sp}$ (non necessarily unique if $k$ is not algebraically closed). This leads us to consider 
\[
[Q_4,\mathrm{Sp}]_{\aone}=\mathrm{KSp}_1(Q_4)=\mathrm{GW}^2_1(Q_4).
\]
We can use the coniveau spectral sequence \cite[Definition 26]{Fasel09c} to compute this group, and we see that (using the fact that the cohomology groups of $Q_4$ with coefficients in strictly $\A^1$-invariant sheaves is easy to compute, see e.g. \cite[Proposition 1.1.5]{Asok16}) this spectral sequence yields a presentation of the form
\[
\mathrm{H}^0(Q_4,\mathbf{GW}_2^2)\to \mathrm{H}^2(Q_4,\mathbf{GW}_3^2)\to \mathrm{GW}_1^2(Q_4)\to 0
\] 
which reduces to 
\[
\KMW2(k)=\mathbf{GW}_2^2(k)\to \mathbf{GW}_1^0(k)\to \mathrm{GW}_1^2(Q_4)\to 0
\] 
If $k$ is algebraically closed, the left-hand term is isomorphic to the uniquely divisible group $\mathbf{K}^{\mathrm{M}}_2(k)$, while the group $ \mathbf{GW}_1^0(k)=\Z/2\Z$ showing that $\mathrm{GW}_1^2(Q_4)=\Z/2\Z$. Since $\Z/2\Z$ is invariant under field extensions, this shows that over any field the direct factor $\Z/2\Z$ of $\mathbf{GW}_1^0(k)$ survives in $\mathrm{GW}_1^2(Q_4)$. 


To describe the image of $\Z/2\Z$ under the map $\mathbf{GW}_1^0(k)\to \mathrm{GW}_1^2(Q_4)$, we draw inspiration from \cite[Section 15.2]{Fasel08a}. The idempotent matrix
\[
M:=\begin{pmatrix} z & 0 & x_2 & -x_1 \\ 0 & z & y_1 & y_2 \\
y_2 & x_1 & 1-z & 0 \\ -y_1 & x_2 & 0 & 1-z\end{pmatrix}
\]
defined in Section \ref{sec:universal2bundle} above yields a rank $2$ projective module $P:=\mathrm{Im}(M)\subset k[Q_4]^4$, and we define a section $s\colon P\to k[Q_4]$ by considering the restriction to $P$ of the projection on the first factor $k[Q_4]^4\to k[Q_4]$. The image of $s$ is generated by the first row of $M$, and is thus the ideal $\mathfrak p:=\langle x_1,x_2,z\rangle$, which is prime of height $2$ with residue field $k(\mathfrak p)=k(y_1,y_2)$.  We have a commutative diagram
\[
\xymatrix{0\ar[r] & k[Q_4]\ar[r]\ar[d]_{-1} & P\ar[r]^-s\ar[d]^-\psi & k[Q_4]\ar[r]\ar@{=}[d] & k[Q_4]/\mathfrak p\ar[r]\ar[d]^-\varphi & 0 \\
0\ar[r] & k[Q_4]\ar[r] & P^\vee\ar[r]^-s & k[Q_4]\ar[r] & \mathrm{Ext}^2_{k[Q_4]}(k[Q_4]/\mathfrak p,k[Q_4])\ar[r] & 0}
\]
in which $\psi$ is the symplectic form induced by $(k[Q_4]^4,H^2)$ on $P$ and $\varphi$ is the symmetric form defined by mapping $1\in k[Q_4]/\mathfrak p$ to the top exact sequence (seen as an element in $ \mathrm{Ext}^2_{k[Q_4]}(k[Q_4]/\mathfrak p,k[Q_4])$). This symmetric form is actually a generator of $\widetilde{\mathrm{CH}}^2(Q_4)$ as shown in \cite[Lemma 4.2.6]{Asok14}.  The permutation matrix $\begin{pmatrix} 0 & 1 \\ 1 & 0\end{pmatrix}$ is orthogonal on $((k[Q_4]/\mathfrak p)^2,\varphi^{\oplus 2})$ and it is straightforward to check that it induces a commutative diagram on the direct sum of the above sequence with itself, corresponding to the permutation of the two factor of $P\oplus P$. This is obviously a symplectic transformation of $(P\oplus P,\psi\oplus\psi)$ and we can use it to define a symplectic matrix $U$ on $k[Q_4]^8$ by the commutative diagram
\[
\xymatrix{k[Q_4]^4\oplus k[Q_4]^4\ar[r]\ar[d]^-U & (P\oplus Q)\oplus (P\oplus Q)\ar[r] & (P\oplus P)\oplus (Q\oplus Q)\ar[d]^-{\tiny{\begin{pmatrix} 0 & 1 & 0 & 0 \\ 1 & 0 & 0 & 0 \\ 0 & 0 & 1 & 0\\ 0 & 0 & 0 & 1\end{pmatrix}}} \\
k[Q_4]^4\oplus k[Q_4]^4\ar[r] & (P\oplus Q)\oplus (P\oplus Q)\ar[r] & (P\oplus P)\oplus (Q\oplus Q)} 
\]
in which $Q$ is the orthogonal of $P$, i.e. the image of the idempotent matrix $N$ above.
To write $U$ explicitly, we observe that the isomorphism $k[Q_4]^4\to P\oplus Q$ is given by $x\mapsto (M(x),N(x))$, with inverse $P\oplus Q\to k[Q_4]^4$ induced by the respective inclusions. Thus
\[
U=\begin{pmatrix} 1-z & 0 & -x_2 & x_1 & z & 0 & x_2 & -x_1 \\ 0 & 1-z & -y_1 & -y_2 & 0 & z & y_1 & y_2 \\
-y_2 & -x_1 & z & 0 & y_2 & x_1 & 1-z & 0 \\ y_1 & -x_2 & 0 & z & -y_1 & x_2 & 0 & 1-z \\
z & 0 & x_2 & -x_1 & 1-z & 0 & -x_2 & x_1 \\ 0 & z & y_1 & y_2 & 0 & 1-z & -y_1 & -y_2 \\
y_2 & x_1 & 1-z & 0 & -y_2 & -x_1 & z & 0 \\ -y_1 & x_2 & 0 & 1-z & y_1 & -x_2 & 0 & z\end{pmatrix}.
\]

We will now reduce the matrix $U\in\mathrm{Sp}_8$ from the previous section into an element in $\mathrm{Sp}_2$ by using elementary symplectic operations. To make the following computations easier, we will need to switch to a different symplectic basis. Let $I_n$ be the $(n \times n)$-identity matrix, we define $J_n := \begin{pmatrix}
    0 & I_n \\ -I_n & 0
\end{pmatrix}$. Setting 
 \begin{equation*}
     S := \begin{pmatrix}
    1& 0& 0& 0& 0& 0& 0& 0\\ 0& 0& 1& 0& 0& 0& 0& 0\\ 0& 0& 0& 0& 1& 0& 0& 0\\ 0& 0& 0& 0& 0& 0& 1& 0\\ 0&
      1& 0& 0& 0& 0& 0& 0\\ 0& 0& 0& 1& 0& 0& 0& 0\\ 0& 0& 0& 0& 0& 1& 0& 0\\ 0& 0& 0& 0& 0& 0& 0& 1
\end{pmatrix},
\end{equation*}
we obtain 
$S \cdot \mathrm{diag(H,H,H,H)} \cdot S^T = J_4$, allowing us to change between the two symplectic bases. 

The elementary symplectic matrices with respect to $J_n$ come in two types. If $E$ is an $(n\times n)$-elementary matrix and $C$ is a \emph{symmetric} $(n\times n)$-matrix. The block matrices 
\begin{equation*}
    \begin{pmatrix}
        E & 0 \\ 0 & E^{-T}
    \end{pmatrix} \quad \text{and} \quad \begin{pmatrix}
        I_n & C \\ 0  & I_n
    \end{pmatrix}
\end{equation*} 
and their transposes are elementary symplectic matrices. On the other hand, the embedding $\mathrm{Sp}_{2n}\to \mathrm{Sp}_{2n+2}$ with respect to the inclusion $J_n \to J_{n+1}$ can be described as follows. The inclusion of the matrix $\begin{pmatrix}
    A & B \\ C  & D
\end{pmatrix} \in Sp_{2n}$ into $Sp_{2n+2}$ is the matrix 
\begin{equation*}
    \begin{pmatrix}
    A & 0_{n \times 1} & B  & 0_{n \times 1} \\ 0_{1 \times n} & 1 & 0_{1 \times n} & 0 \\ C & 0_{n \times 1}  & D & 0_{n \times 1} \\ 0_{1 \times n} & 0 & 0_{1 \times n} & 1
\end{pmatrix}.
\end{equation*}

We will now start with the symplectic reduction of the matrix
\begin{equation*}
    U=\begin{pmatrix} 1-z & 0 & -x_2 & x_1 & z & 0 & x_2 & -x_1 \\0 & 1-z & -y_1 & -y_2 & 0 & z & y_1 & y_2 \\
-y_2 & -x_1 & z & 0 & y_2 & x_1 & 1-z & 0 \\ y_1 & -x_2 & 0 & z & -y_1 & x_2 & 0 & 1-z \\
z & 0 & x_2 & -x_1 & 1-z & 0 & -x_2 & x_1 \\0 & z & y_1 & y_2 & 0 & 1-z & -y_1 & -y_2 \\
y_2 & x_1 & 1-z & 0 & -y_2 & -x_1 & z & 0 \\ -y_1 & x_2 & 0 & 1-z & y_1 & -x_2 & 0 & z\end{pmatrix}.
\end{equation*}
Performing the above change of basis, we see that $U$ becomes symplectic with respect to $J_4$ by conjugating with $S$:

\begin{equation*}
    M_1 := SUS^T = \begin{pmatrix}
        
    -z+1& -x_2& z& x_2& 0& x_1& 0& -x_1\\ -y_2& z& y_2& -z+1& -x_1& 0& x_1& 0\\ z& x_2& -z+1& -x_2& 0& -x_1& 0& x_1\\ y_2& -z+1& -y_2& z& x_1& 0& -x_1& 0\\ 0& -y_1& 0& y_1& -z+1& -y_2& z& y_2\\ y_1& 0& -y_1& 0& -x_2& z& x_2& -z+1\\ 0& y_1& 0& -y_1& z& y_2& -z+1& -y_2\\ -y_1& 0& y_1& 0& x_2& -z+1& -x_2& z \end{pmatrix}.
\end{equation*}
Setting
\begin{equation*}
    E_1 := \begin{pmatrix}
        1&0&0&0\\0&1&0&1\\0&0&1&0\\0&0&0&1
    \end{pmatrix}, \quad E_2 := \begin{pmatrix}
        1& 0& 0& 0\\ 0& 1& 0& 0\\ 0& 0& 1& 0\\ -y_2& z-1& y_2& 1
    \end{pmatrix},
\end{equation*}
the product 
\begin{equation*}
    \begin{pmatrix}
        E_1 & 0 \\ 0 & E_1^{-T}
    \end{pmatrix}^{-1}M\begin{pmatrix}
        E_1 & 0 \\ 0 & E_1^{-T}
    \end{pmatrix} \begin{pmatrix}
        E_2 & 0 \\ 0 & E_2^{-T} 
    \end{pmatrix}
\end{equation*}
yields the matrix 
\begin{equation*}
    \begin{pmatrix}
        -z+1& -x_2& z& 0& 0& 2x_1& 0& -2x_1z+x_1\\ -2y_2& 2z-1& 2y_2& 0& -2x_1& 0& 2x_1& -4x_1y_2\\ z& x_2& -z+1& 0& 0& -2x_1& 0& 2x_1z-x_1\\ 0& 0& 0& 1& x_1& 0& -x_1&2x_1y_2\\ 0& -y_1& 0& 0& -z+1& -2y_2& z& 0\\ y_1& 0& -y_1& 0& -x_2& 2z-1& x_2& -2x_2y_2-2z^2+2z\\ 0& y_1& 0& 0& z& 2y_2& -z+1& 0\\ 0& 0& 0& 0& 0& 0& 0& 1

.
    \end{pmatrix}.
\end{equation*}
multiplying from the right by $\begin{pmatrix}
    I_4 & S_1 \\ 0 & I_4
\end{pmatrix}$, where $S_4$ is the symmetric matrix 
\begin{equation*}
    S_4 = \begin{pmatrix}
        0&0&0&-x_1\\0&0&0&0\\0&0&0&x_1\\-x_1&0&x_1&-2x_1y_2
    \end{pmatrix}
\end{equation*}
we obtain 
\begin{equation*}
    \begin{pmatrix}
        -z+1& -x_2& z& 0& 0& 2x_1& 0& 0\\ -2y_2& 2z-1& 2y_2& 0& -2x_1& 0& 2x_1& 0\\ z& x_2& -z+1& 0& 0& -2x_1& 0& 0\\ 0& 0& 0& 1& 0& 0& 0& 0\\ 0& -y_1& 0& 0& -z+1& -2y_2&z& 0\\ y_1& 0& -y_1& 0& -x_2& 2z-1& x_2& 0\\ 0& y_1& 0& 0& z& 2y_2& -z+1& 0\\ 0& 0& 0& 0& 0& 0& 0& 1        
    \end{pmatrix},
\end{equation*}
from which we extract the $\mathrm{Sp}_6$ matrix 
\begin{equation*}
    M_2 := \begin{pmatrix}
        -z+1& -x_2& z& 0& 2x_1& 0\\ -2y_2& 2z-1& 2y_2& -2x_1& 0& 2x_1\\ z& x_2& -z+1& 0& -2x_1& 0\\ 0& -y_1& 0& -z+1& -2y_2& z\\ y_1& 0& -y_1& -x_2& 2z-1& x_2\\ 0& y_1& 0& z&2y_2& -z+1
    \end{pmatrix}.
\end{equation*}
Repeating the same procedure on $M_2$ with the elementary matrices 
\begin{equation*}
    E_3 := \begin{pmatrix}
        1 & 0 & 1 \\ 0 & 1 &0 \\ 0 & 0 & 1
    \end{pmatrix}, \quad E_4 := \begin{pmatrix}
        1 & 0 & 0 \\ 0 & 1 & 0 \\ -z & -x_2 & 1
    \end{pmatrix}
\end{equation*}
and the symmetric matrix
\begin{equation*}
   S_2 :=  \begin{pmatrix}
        0& 0& 0\\ 0& 0& 2x_1\\ 0& 2x_1& 2x_1x_2
    \end{pmatrix}
\end{equation*}
yields the $\mathrm{Sp}_4$ matrix
\begin{equation*}
    M_3 := \begin{pmatrix}
    -2z+1& -2x_2& 0& 4x_1\\ -2y_2& 2z-1& -4x_1& 0\\ 0& -y_1& -2z+1& -2y_2\\ y_1& 0& -2x_2& 2z-1
\end{pmatrix}.
\end{equation*}

Consider next the elementary matrices 
\begin{equation*}
    E_5 := \begin{pmatrix}
        1 & 2x_2 \\0 & 1 
    \end{pmatrix}, \quad E_6 := \begin{pmatrix} 1 & 0 \\ -y_1 & 1
    \end{pmatrix} 
\end{equation*}
and the symmetric matrices 
\begin{equation*}
    S_3 := \begin{pmatrix}
        0 & 0 \\ 0 & -1 
    \end{pmatrix}, \quad S_4 := \begin{pmatrix}
        0 & 0 \\ 0 & 1 - (2z-1).
    \end{pmatrix}
\end{equation*}
The product
\begin{equation*} M_4 := 
\begin{pmatrix} 
        E_5 & 0 \\ 0 & E_5^{-T}
    \end{pmatrix}
    \begin{pmatrix}
        I_2 & 0 \\ S_3 & I_2 
    \end{pmatrix} M_3 \begin{pmatrix}
        I_2 & S_4 \\ 0 & I_2 
    \end{pmatrix} \begin{pmatrix}
        E_6 & 0 \\ 0 & E_6^{-T}
    \end{pmatrix} 
\end{equation*}
yields the matrix 
\begin{equation*}
    M_4 = \begin{pmatrix}
       -4x_2y_1z+4x_2y_1-4x_2y_2-2z+1& 4x_2z-4x_2& -8x_1x_2& 8x_2^2y_2+8x_2z+4x_1-8x_2\\ -2y_1z+y_1-2y_2& 2z-1& -4x_1& 4x_2y_2+2z-2\\ y_1^2& -y_1& -2z+1& -y_1-2y_2\\-2x_2y_1^2+2y_1z+2y_2& 2x_2y_1-2z+1& 4x_2z+4x_1-4x_2& 1
    \end{pmatrix}.
\end{equation*}
Using the $1$ in the bottom right corner, we aim to use elementary transformation to remove the entries $-y_1 - 2y_2$ and $4x_2z + 4x_1 -4x_2$. Consider next the elementary matrices 
\begin{equation*}
    E_7 := \begin{pmatrix}
        1  & 2x_2 \\ 0 & 1 
    \end{pmatrix}, \quad E_8 := \begin{pmatrix}
        1 & 0 \\ -y_1 & 1
    \end{pmatrix}
\end{equation*}
we set 
\begin{equation*}
    N := \begin{pmatrix}
        E_7 & 0 \\ 0 & E_7^{-T}
    \end{pmatrix}
    M_4\begin{pmatrix}
        E_8& 0 \\ 0 & E_8^{-T}
    \end{pmatrix} .
\end{equation*}
The matrix $N$ is too big to write down explicitly here, but it can be written as the matrix
\begin{equation*}
    \begin{pmatrix}
        & & & N_{14} \\
        & N'& & N_{24} \\
        & & & 0 \\
        N_{41} & N_{42} & 0 & 1
    \end{pmatrix}
\end{equation*}
where $N'$ is a $(3 \times 3)$-matrix. 
Using finally the symmetric matrices 
\begin{equation*}
    S_5 = \begin{pmatrix}
        0 & -N_{14} \\ -  N_{14} & -N_{24}
    \end{pmatrix},\quad S_6 = \begin{pmatrix}
        0 & -N_{41} \\ -  N_{41} & -N_{42}
    \end{pmatrix} .
\end{equation*} 
%\jean{Are you sure about the matrix $S_5$? Why is it symmetric?}
The product 
\begin{equation*}
    \begin{pmatrix}
        I_2 & S_5 \\ 0 & I_2
    \end{pmatrix}N \begin{pmatrix}
        I_2 & 0 \\ S_6 & I_2
    \end{pmatrix}
\end{equation*} 
is then the image of the $Sp_2 $ matrix $\begin{pmatrix}
    a & b \\ c & d
\end{pmatrix}$ into $Sp_4$ where 

\begin{align*}
    a =& 16x_2^3y_1^2y_2+16x_2^2y_1^2z-16x_2^2y_1y_2z-16x_2^2y_1^2-8x_2^2y_1y_2-16x_2^2y_2^2-24x_2y_1z^2 \\&+20x_2y_1z-8x_2y_2z+8z^3+4x_2y_1-8x_1y_2+12x_2y_2-8z^2-2z+1, \\
    b =& -32x_2^3y_2z-32x_1x_2^2y_2+32x_2^3y_2-32x_2^2z^2-48x_1x_2z+64x_2^2z-16x_1^2+40x_1x_2-32x_2^2, \\
    c =& -2x_2y_1^3-4x_2y_1^2y_2+2y_1^2z+4y_1y_2z+y_1^2+2y_1y_2+4y_2^2, \\
    d =& 4x_2y_1z+8x_2y_2z-4x_2y_1+8x_1y_2-12x_2y_2-4z^2+2z+1.
\end{align*}

\subsection{Second construction: weight shifting}


\subsubsection{The universal bundle on $Q_5$}
We consider the fiber sequence
\[
Q_5\to \mathrm{BSL}_2\to \mathrm{BSL}_3 
\]
in which the left-hand map classifies the universal rank $2$ bundle $P$ on $Q_5$. Writing 
\[
R=k[x_1,x_2,x_3,y_1,y_2,y_3]/(\sum_{i=1}^3 x_iy_i-1)
\]
for the coordinate ring of $Q_5$, it is given by the exact sequence
\[
0\to P\to R^3\xrightarrow{(x_1\phantom{i}x_2\phantom{i}x_3)} R\to 0,
\]
which is split by the homomorphism $R\xrightarrow{\begin{pmatrix} y_1 \\ y_2 \\y_3\end{pmatrix}}R^3$. We can then see $P$ as the image of the idempotent matrix
\begin{equation}\label{eqn:universalidempotent}
M(P):=\begin{pmatrix} 1-x_1y_1 & -x_2y_1 & -x_3y_1 \\ -x_1y_2 & 1-x_2y_2 & -x_3y_2 \\ -x_1y_3 & -x_2y_3 & 1-x_3y_3\end{pmatrix}.
\end{equation}
\begin{rem}
It is possible to show that the map $Q_5\to \mathrm{BSL}_2$ factorizes as $Q_5\to Q_4\to \mathrm{BSL}_2$, in which $Q_5\to Q_4$ is the suspension of the exotic Hopf map constructed in the previous section and $Q_4\to \mathrm{BSL}_2$ is the map classifying the universal rank $2$ symplectic bundle on $Q_4$. Since we don't use this fact, we leave the details to the reader.
\end{rem}

The scheme $\A^3\smallsetminus\{0\}$ can be covered by the two open subschemes $\A^2\smallsetminus\{0\}\times \A^1$ and $\A^2\times \gm{}$, and we transform this covering into a covering of $Q_5$ using the affine open subscheme $U:=\{x_3\neq 0\}\subset Q_5$ and the affine closed subscheme $V:=\{y_3=0\}\subset Q_5$. We obtain a Cartesian square 
\begin{equation}\label{eqn:homotopycartesianQ5}
\xymatrix{U\cap V\ar[r]\ar[d] & U\ar[d]^-j \\
V\ar[r]_-i & Q_5}
\end{equation}
in which $U\cap V$ is the affine scheme $Q_3\times \gm$, with coordinates $(x_1,x_2,y_1,y_2,x_3)$. The projections $U\to \A^2\times \gm{}$ and $V\to \A^2\smallsetminus\{0\}\times \A^1$ induced by the projections $Q_5\to \A^3\smallsetminus\{0\}$ are weak equivalences, and it follows that the maps $i$ and $j$ are homotopy trivial. Consequently, the pull-back of the universal rank $2$ bundle on $Q_5$ along both $i$ and $j$ is trivial and we now find explicit trivializations starting with $U$. Consider the matrix
\[
E:=\begin{pmatrix} 0 & 0 & -x_3 \\ 0 & 1 & 0 \\ x_3^{-1} & -x_2x_3^{-1} & x_1\end{pmatrix}\begin{pmatrix} 1 & 0 & 0 \\ y_2 & 1 & 0 \\ -x_3^{-1}y_1 & 0 & 1\end{pmatrix}=\begin{pmatrix} y_1 & 0 & -x_3 \\ y_2 & 1 & 0 \\ y_3 & -x_2x_3^{-1} & x_1 \end{pmatrix}
\]
which is easily seen to be elementary (using Whitehead Lemma). We have 
\[
E^{-1}=\begin{pmatrix} x_1 & x_ 2 & x_3 \\ -x_1y_2 & 1-x_2y_2 & -x_3y_2 \\ -x_3^{-1}(1-x_1y_1) & x_3^{-1}x_2y_1 & y_1\end{pmatrix}
\]
and the properties $(x_1\phantom{i}x_2\phantom{i}x_3)E=(1\phantom{i}0\phantom{i}0)$, $E^{-1}\begin{pmatrix} y_1 \\ y_2 \\y_3\end{pmatrix}=\begin{pmatrix} 1 \\ 0 \\0\end{pmatrix}$ are easily verified. It follows that $E$ induces an isomorphism $k[U]^2\to P_{\vert U}$, for instance since
\[
E^{-1}M(P)E=E^{-1}(\mathrm{Id}-\begin{pmatrix} y_1 \\ y_2 \\y_3\end{pmatrix}(x_1\phantom{i}x_2\phantom{i}x_3))E=\mathrm{Id}-\begin{pmatrix} 1 \\ 0 \\0\end{pmatrix}(1\phantom{i}0\phantom{i}0)=\begin{pmatrix} 0 & 0 & 0 \\ 0 & 1 & 0 \\ 0 & 0 & 1\end{pmatrix}.
\]
We perform the same task over $V$, now using the matrix
\begin{eqnarray*}
F& := & \begin{pmatrix}1 & 0 & (1-x_3)y_1 \\ 0 & 1 & (1-x_3)y_2 \\ 0 & 0 & 1\end{pmatrix}\begin{pmatrix}1 & 0 & 0 \\ 0 & 1 & 0 \\ -x_1 & -x_2 & 1\end{pmatrix}\begin{pmatrix}0 & 0 & -1 \\ 0 & 1 & 0 \\ 1 & 0 & 0\end{pmatrix}\begin{pmatrix} 1 & 0 & 0 \\ y_2 & 1 & 0 \\ -y_1 & 0 & 1\end{pmatrix}\\
& = & \begin{pmatrix} y_1 & -x_2y_1(1-x_3) & -1+(1-x_3)x_1y_1 \\ y_2 & 1-x_2y_2(1-x_3) & (1-x_3)x_1y_2 \\ 0 & -x_2 & x_1\end{pmatrix}
\end{eqnarray*}
which is also elementary. We find
\[
F^{-1}=\begin{pmatrix} x_1 & x_ 2 & x_3 \\ -x_1y_2 & x_1y_1 & -y_2 \\ -x_2y_2 & x_2y_1 & y_1\end{pmatrix}
\]
and $(x_1\phantom{i}x_2\phantom{i}x_3)F=(1\phantom{i}0\phantom{i}0)$, $F^{-1}\begin{pmatrix} y_1 \\ y_2 \\y_3\end{pmatrix}=F^{-1}\begin{pmatrix} y_1 \\ y_2 \\0\end{pmatrix}=\begin{pmatrix} 1 \\ 0 \\0\end{pmatrix}$ in the algebra 
\[
k[V]=k[x_1,x_2,y_1,y_2,x_3]/(x_1y_1+x_2y_2-1).
\]
As above, $F$ induces an isomorphism $k[V]^2\to P_{\vert V}$ since $F^{-1}M(P)F=\begin{pmatrix} 0 & 0 & 0\\ 0 & 1 & 0 \\ 0 & 0 & 1\end{pmatrix}$.

Over the intersection $U\cap V$, $E$ and $F$ do the same job and the above equations show that $E^{-1}F$ induces a matrix in $\mathrm{SL}_2(U\cap V)$. Explicitly, we find
\[
E^{-1}F=\begin{pmatrix} 1 & 0 & 0 \\ 0 &  1-x_2y_2(1-x_3) & x_1y_2(1-x_3) \\
0 & (x_3^{-1}-1)x_2y_1 & x_3^{-1}-x_1y_1(x_3^{-1}-1)\end{pmatrix}
\]
and we set 
\[
T:=\begin{pmatrix} 1-x_2y_2(1-x_3) & x_1y_2(1-x_3) \\ (x_3^{-1}-1)x_2y_1 & x_3^{-1}-x_1y_1(x_3^{-1}-1)\end{pmatrix}\in \mathrm{SL}_2(U\cap V).
\]
A simple computation yields 
\[
T:=\begin{pmatrix} 1 & 0 \\ 0 & x_3^{-1}\end{pmatrix}\begin{pmatrix} x_1 & -y_2 \\ x_2 & y_1\end{pmatrix}\begin{pmatrix} 1 & 0 \\ 0 & x_3^{-1}\end{pmatrix}^{-1}\begin{pmatrix} x_1 & -y_2 \\ x_2 & y_1\end{pmatrix}^{-1},
\]
i.e. $T$ is a commutator.
The base point $1\in\gm{}$ yields an embedding $Q_3\to Q_3\times \gm{}$, while the base point $(1,0,1,0)\in Q_3$ induces an embedding $\gm{}\to Q_3\times \gm{}$. A direct computation shows that the restriction of $T$ along both these morphisms is trivial. Since $\mathrm{SL}_2$ is a sheaf, the morphism $T\colon Q_3\times \gm{}\to \mathrm{SL}_2$ then induces a morphism of sheaves
\[
T\colon Q_3\wedge \gm{}\to \mathrm{SL}_2
\]
that we still denote by $T$, slightly abusing notation.  
Now, the square \eqref{eqn:homotopycartesianQ5} is homotopy cocartesian, being weakly equivalent to the cocartesian square 
\[
\xymatrix{(\A^2\smallsetminus \{0\})\times \gm{}\ar[r]\ar[d] & \A^2\times \gm{}\ar[d] \\
(\A^2\smallsetminus \{0\})\times \A^1 \ar[r] & \A^3\smallsetminus\{0\},}
\] 
and the long exact sequence (of pointed sets and groups) associated to this homotopy cocartesian square induces in particular a map
\[
\xymatrix{[(\A^2\smallsetminus \{0\})\times \gm{},\mathrm{SL}_2]_{\A^1}\ar[r] &[\A^3\smallsetminus\{0\},\mathrm{BSL}_2]_{\A^1}.}
\]
The above discussion shows that the image of $T\in [Q_3\wedge \gm{},\mathrm{SL}_2]_{\A^1}$ under this map is precisely the (homotopy class of) the universal rank $2$ bundle on $Q_5$. In other words, $T$ is the adjoint of the morphism in $[Q_5,\mathrm{BSL}_2]_{\A^1}$ that gives the universal rank $2$ bundle. On the other hand, we have
\[
[\A^3\smallsetminus\{0\},\mathrm{BSL}_2]_{\A^1}=\piaone_{2,3}(\mathrm{BSL}_2)(k)=\KMW {-1}(k)
\]
and the class of the universal bundle is a generator of this group. It follows that $h\cdot T$ is trivial, or in other terms $T$ is $2_{\epsilon}$-torsion. 

\subsubsection{Weight shifting}\label{sec:weightshifting}

As seen in the above section, the morphism $T\colon Q_3\wedge \gm{}\to \mathrm{SL}_2$ is $h=2_{\epsilon}$-torsion. Writing $m\colon Q_3\wedge \gm{}\to Q_3\wedge \gm{}$ the morphism $\mathrm{Id}\wedge (\cdot 2)$, this means by \cite[Proposition 2.1.9]{Asok20b} that the composite 
\[
Q_3\wedge \gm{}\xrightarrow{m}Q_3\wedge \gm{}\xrightarrow{T} \mathrm{SL}_2.
\]
is trivial (up to homotopy). It follows that $T$ is in the kernel ${}_{2_\epsilon}\piaone_{1,3}(\mathrm{SL}_2)$ of the map 
\[
\piaone_{1,3}(\mathrm{SL}_2)=[Q_3\wedge \gm{},\mathrm{SL}_2]_{\A^1}\xrightarrow{2_{\epsilon}} [Q_3\wedge \gm{},\mathrm{SL}_2]_{\A^1}=\piaone_{1,3}(\mathrm{SL}_2)
\] 
induced by $m$. As explained in \cite[Theorem 2.1.11]{Asok20b}, the choice of $\tau=-1$ yields a homomorphism
\[
{}_{2_\epsilon}\piaone_{1,3}(\mathrm{SL}_2)\to \piaone_{2,2}(\mathrm{SL}_2)/[-1]\piaone_{2,3}(\mathrm{SL}_2)
\] 
that is called \emph{weight shifting}. Our purpose in this section is to explicitly perform this procedure that is a priori quite abstract.


For this, it is slightly inconvenient to work with $Q_3\wedge \gm{}$ (which is just a sheaf), and we instead compute the composite
\[
Q_3\times \gm{}\xrightarrow{m'}Q_3\times \gm{}\xrightarrow{T} \mathrm{SL}_2
\]
where $m'$ is a model of $m$. To obtain this map, we consider the morphism $m'\colon Q_5\to Q_5$ defined by
\[
m'(x_i)=\begin{cases} x_i & \text{ for $i=1,2$,} \\ x_3^2 & \text{ for $i=3$,}\end{cases}; \phantom{i} m'(y_i)=\begin{cases} y_i(1+x_3y_3) & \text{ if $i=1,2$.} \\ y_3^2 & \text{ if $i=3$.}\end{cases}
\]
and we observe that it respects the subschemes $U$ and $V$, i.e. it induces endomorphisms of $U$ and $V$, and consequently of $U\cap V$. 
We have 
\[
T\circ m'=\begin{pmatrix} 1-x_2y_2(1-x_3^2) & x_1y_2(1-x_3^2) \\ x_2y_1(x_3^{-2}-1) & x_3^{-2}-x_1y_1(x_3^{-2}-1)\end{pmatrix}
\]
that we may rewrite as follows using the above presentation as a commutator
\[
T\circ m'=\begin{pmatrix} x_3 & 0 \\ 0 & x_3^{-1}\end{pmatrix}\begin{pmatrix} x_1 & -y_2 \\ x_2 & y_1\end{pmatrix}\begin{pmatrix} x_3^{-1} & 0 \\ 0 & x_3\end{pmatrix}\begin{pmatrix} x_1 & -y_2 \\ x_2 & y_1\end{pmatrix}^{-1}.
\]
Since $Q_3\times\gm{}$ is affine and $\mathrm{SL}_2$ is $\A^1$-naive (e.g. \cite{Asok15a}), there should be an explicit naive $\A^1$-homotopy between $T\circ m'$ and the identity. We may use Whitehead lemma to obtain
\[
\begin{pmatrix} x_3 & 0 \\ 0 & x_3^{-1}\end{pmatrix}=\begin{pmatrix} 1 & 0 \\ x_3^{-1} & 1\end{pmatrix}\begin{pmatrix} 1 & 1-x_3  \\ 0 & 1\end{pmatrix}\begin{pmatrix} 1 & 0 \\ -1 & 1\end{pmatrix}\begin{pmatrix} 1 & 1-x_3^{-1}  \\ 0 & 1\end{pmatrix},
\]
from which we deduce the expected homotopy. More precisely, we consider the matrix
\[
A(t):=\begin{pmatrix} 1 & 0 \\ tx_3^{-1} & 1\end{pmatrix}\begin{pmatrix} 1 & t(1-x_3)  \\ 0 & 1\end{pmatrix}\begin{pmatrix} 1 & 0 \\ -t & 1\end{pmatrix}\begin{pmatrix} 1 & t(1-x_3^{-1})  \\ 0 & 1\end{pmatrix}, 
\]
and we obtain a matrix
\[
D(t):=A(t)\begin{pmatrix} x_1 & -y_2 \\ x_2 & y_1\end{pmatrix}A(t)^{-1}\begin{pmatrix} x_1 & -y_2 \\ x_2 & y_1\end{pmatrix}^{-1}
\]
which has the property that $D(0)=\mathrm{Id}$ and $D(1)=T\circ m'$. Its restriction to both $Q_3\times \{1\}\times \A^1$ and $\{1,0,1,0\}\times \gm{}\times \A^1$ is trivial, and consequently $D$ can be seen as a map 
\[
(Q_3\wedge \gm{})\times \A^1\to \mathrm{SL}_2.
\] 
Setting $x_3=-1$, we obtain a map $Q_3\to Q_3\wedge \gm{}$, which in turn yields a map $G(t)\colon Q_3\times \A^1\to \mathrm{SL}_2$, which becomes trivial (since $D(1)=T\circ m'$ and $m'(-1)=1$) when restricted to $Q_3\times S^0$, where we set $S^0:=\{t(1-t)=0\}\subset \A^1$. Thus, we obtain a map $(Q_3\times\A^1)/(Q_3\times S^0)\to \mathrm{SL}_2$ that is explicitly of the form
\[
G(t)  =  \begin{pmatrix} 1-2t^2 & 4(t-t^3) \\ -2(t-t^3) & 1-6t^2+4t^4 \end{pmatrix}\begin{pmatrix} x_1 & -y_2 \\ x_2 & y_1\end{pmatrix}\begin{pmatrix} 1-6t^2+4t^4 & 4(t^3-t) \\ 2(t-t^3) & 1-2t^2 \end{pmatrix}\begin{pmatrix} x_1 & -y_2 \\ x_2 & y_1 \end{pmatrix}^{-1}. 
\]
In more traditional terms, $G(t)\in \mathrm{SL}_2(k[Q_3][t],t(1-t))$. Now, remember that we have a morphism
\[
\alpha_3\colon Q_3\times \A^1\to Q_4
\]
given by $\alpha_3(x_1,x_2,y_1,y_2,t)=(x_1,x_2,y_1t(1-t),y_2t(1-t),t)$, mapping $ Q_3\times S^0$ to the closed subset  $\{z(1-z)=0\}\subset Q_4$. On the other hand, the above morphism induces an isomorphism 
\[
\alpha_3\colon Q_3\times (\A^1\smallsetminus \{0,1\})\simeq D(z(1-z))\subset Q_4
\]
with inverse given by $(x_1,x_2,y_1,y_2,z)\mapsto (x_1,x_2,\frac{y_1}{z(1-z)},\frac{y_2}{z(1-z)},z)$. The matrix $G(t)$ restricts to a map
\[
Q_3\times (\A^1\smallsetminus \{0,1\})\to \mathrm{SL}_2
\]
and thus induces a map $D(z(1-z))\to \mathrm{SL}_2$ that should morally be extended to $Q_4$ by setting its restriction to $\{z(1-z)=0\}$ to be trivial. However, $G(t)$ lifts to a map on $Q_4$ only if every monomial involving $\frac{y_i}{z(1-z)}$ is multiplied by $z(1-z)$, which is not the case as a simple computation shows. We can nevertheless modify $G(t)$ up to homotopy relative to $t(1-t)$ to achieve this, using the next lemma.  

\begin{lem}
Let $M:=\begin{pmatrix} x_1 & -y_2 \\ x_2 & y_1\end{pmatrix}$. Then:
\[
M\begin{pmatrix} s & p \\ q & r\end{pmatrix}M^{-1}=s\begin{pmatrix} x_1 \\ x_2 \end{pmatrix}\begin{pmatrix} y_1 & y_2\end{pmatrix}+ p\begin{pmatrix} x_1 \\ x_2 \end{pmatrix}\begin{pmatrix} -x_2 & x_1\end{pmatrix} + q\begin{pmatrix} -y_2 \\ y_1 \end{pmatrix}\begin{pmatrix} y_1 & y_2\end{pmatrix} + r\begin{pmatrix} -y_2 \\ y_1 \end{pmatrix}\begin{pmatrix} -x_2 & x_1\end{pmatrix}
\] 
for any $s,p,q,r\in k[t]$.
\end{lem}

Now, 
\[
G(t)=\begin{pmatrix} 1-2t^2 & 4(t-t^3) \\ -2(t-t^3) & 1-6t^2+4t^4 \end{pmatrix}\begin{pmatrix} x_1 & -y_2 \\ x_2 & y_1\end{pmatrix}\begin{pmatrix} 1-6t^2+4t^4 & 4(t^3-t) \\ 2(t-t^3) & 1-2t^2 \end{pmatrix}\begin{pmatrix} x_1 & -y_2 \\ x_2 & y_1 \end{pmatrix}^{-1}.
\]
Using this lemma, we replace $\begin{pmatrix} 1-6t^2+4t^4 & 4(t^3-t) \\ 2(t-t^3) & 1-2t^2 \end{pmatrix}$ by a matrix $B(t) = \begin{pmatrix} s & p \\ q & r\end{pmatrix}$ for some polynomials $s,p,q,r \in k[t]$ assumed to be invertible and homotopic to the latter relative to $t(1-t)$, and having the properties that the product 
\[
\begin{pmatrix} x_1 & -y_2 \\ x_2 & y_1\end{pmatrix}B(t)\begin{pmatrix} x_1 & -y_2 \\ x_2 & y_1\end{pmatrix}^{-1}
\] 
lies in the image of the $k$-algebra map $\psi \colon k[Q_4] \to k[Q_3\times \A^1]$ induced by $\alpha_3$, i.e.  
\[
x_i \mapsto x_i, y_i \mapsto t(1-t)y_i, z \mapsto t.
\]
Let then 
\[ 
B(t) = \begin{pmatrix} s & p \\ q & r\end{pmatrix} = \left (
\begin{pmatrix} 1 & \alpha(1-t)-2t \\ 0 & 1\end{pmatrix}\begin{pmatrix} 1 & 0 \\ t+\beta t(1-t) & 1\end{pmatrix}\right)^2
\] 
for some polynomials $\alpha,\beta\in k[t]$. A direct computation yields 
\[
q= (2(t+t(1-t)\beta)+(t+t(1-t)\alpha)^2(t(1-t)\alpha-2t)).
\] 
If we want the expression to be divisible by $t^2(1-t)^2$, we may for instance take $\alpha=-6$, $\beta=-1$ to obtain a matrix
\[
B(t) := \left(\begin{pmatrix}
    1 & -2t - 6t(1-t) \\ 0 & 1 
\end{pmatrix}\begin{pmatrix}
    1 & 0 \\ t - t(1-t) & 1
\end{pmatrix}\right)^2 = \left(\begin{pmatrix}
    1 & 6t^2-8t \\ 0 & 1 
\end{pmatrix}\begin{pmatrix}
    1 & 0 \\ t^2 & 1
\end{pmatrix}\right)^2. 
\]
that is apparently homotopic to $A(t)^{-1}$ relative to $t(1-t)$. We can then set 
\[
G'(t) = \begin{pmatrix} 1-2t^2 & 4(t-t^3) \\ -2(t-t^3) & 1-6t^2+4t^4 \end{pmatrix}\begin{pmatrix} x_1 & -y_2 \\ x_2 & y_1\end{pmatrix}B(t)\begin{pmatrix} x_1 & -y_2 \\ x_2 & y_1\end{pmatrix}^{-1}
\] 
and observe that this matrix indeed extends to a map $Q_4\to  \mathrm{SL}_2$ as its coordinates lie in the image of the $k$-algebra map $\psi \colon k[Q_4] \to k[Q_{3}\times \A^1]$. By computing $\psi^{-1}(G'(t))$ we get a map $Q_4 \to Q_3$ given by the unimodular row $(a,b)$, where
\begin{align*}
   a =& 144x_2^2z^9+144x_2y_1z^9+72x_1x_2z^8-384x_2^2z^8-240x_2y_1z^8-72x_2 y_2z^8-72z^{10}-192x_1x_2z^7\\ &+112x_2^2z^7-128x_2y_1z^7+120x_2y_2z^7+192 z^9+92x_1x_2z^6+384x_2^2z^6+256x_2y_1z^6+28x_2y_2z^6\\&-92z^8+96x_1x_2z ^5-208x_2^2z^5+48x_2y_1z^5-24y_1^2z^5-68x_2y_2z^5-96z^7-40x_1x_2z^4- 64x_2^2z^4\\&-16x_2y_1z^4-16y_1^2z^4-28x_2y_2z^4+12y_1y_2z^4+28z^6-32 x_1x_2z^3-48x_2^2z^3-64x_2y_1z^3+16y_1^2z^3\\&+4x_2y_2z^3+8y_1y_2z^3+48 z^5-12x_1x_2z^2+64x_2^2z^2+16y_1^2z^2+16x_2y_2z^2-2y_1y_2z^2+18z^4\\&+ 16x_1x_2z+8y_1^2z-4y_1y_2z-24z^3-2y_1y_2-2z^2+1, \\
   b = & -144x_1x_2z^9+144x_2y_2z^9-72x_1^2z^8+384x_1x_2z^8+72x_1y_2z^8-240x_2 y_2z^8+192x_1^2z^7- 112x_1x_2z^7\\&-120x_1y_2z^7-128x_2y_2z^7-92x_1^2z^6 -384x_1x_2z^6-28x_1y_2z^6+256x_2y_2z^6-96x_1^2z^5+208x_1x_2z^5\\&+68x_1 y_2z^5+48x_2y_2z^5-24y_1y_2z^5-24z^7+40x_1^2z^4+64x_1x_2z^4+28x_1y_2 z^4-16x_2y_2z^4\\&-16y_1y_2z^4+12y_2^2z^4+32z^6+32x_1^2z^3+48x_1x_2z^3- 4x_1y_2z^3-64x_2y_2z^3+16y_1y_2z^3+8y_2^2z^3\\&+24z^5+12x_1^2z^2-64x_1 x_2z^2-16x_1y_2z^2+16y_1y_2z^2-2y_2^2z^2-32z^4-16x_1^2z+8y_1y_2z\\&-4 y_2^2z-4z^3-2y_2^2+4z.
\end{align*}

\begin{rem}
We note that both the exotic Hopf maps obtained in Section \ref{sec:symplecticHopfmap} and above are defined over $\Z$. However,  they reduce to the trivial map on $\mathbb{F}_2$ as a direct check shows. In the second case, this can be explained by the fact that we had to choose a primitive $2$nd root of $1$. In the first case, we used a computation of $\mathrm{GW}_1^0(k)$ that needs $2$ to be invertible, and the generator we used is trivial in $\mathrm{GW}_1^0(\mathbb{F}_2)$.
\end{rem}

\begin{rem}
A direct computation shows that the above map is not trivial relative $z(1-z)$. This could probably be further achieved by modifying it up to homotopy, but we haven't carried out the necessary computations since we don't need this property in the sequel.
\end{rem}

\begin{rem}
It would probably be difficult to prove that the exotic Hopf maps of Section \ref{sec:symplecticHopfmap} and the one above are homotopic. They however coincide after base change to an algebraically closed field as proved below.
\end{rem}

%\jean{We may observe several things:
%\begin{enumerate}
%\item The map is non trivial relative to $z(1-z)$. This could probably be further achieved by modifying it up to homotopy. For instance, for $z=0$, we obtain a row $(1-2y_1y_2,2y_2^2)$, which is the first row of the matrix
%\[
%\begin{pmatrix}1-2y_1y_2 & 2y_2^2 \\ -2y_1^2 & 1+2y_1y_2 \end{pmatrix}
%\] 
%which is homotopically trivial.
%\item It is not clear if it is homotopic to the other exotic Hopf maps described in this paper (including Turiel's).  We may still say that over an algebraically closed field they coincide since they both correspond to the only non trivial element.
%\end{enumerate} }


\section{A rank $2$ bundle on $\Proj^3_k$}

We finally put everything together to construct the rank $2$ bundle on $\mathbb{J}^3$. Indeed, we consider the following composite
\[
\mathbb{J}^3\xrightarrow{\pi_3} Q_6\xrightarrow{\Sigma_{\pone}\eta'}Q_5\to \mathrm{BSL}_2
\]
where $\pi_3$ is the morphism described in Section \ref{sec:collapsemap}, and $\Sigma_{\pone}\eta'$ is the $\pone$-suspension of either the exotic Hopf map of Section \ref{sec:weightshifting} or Section \ref{sec:symplecticHopfmap}. Explicitly, it is given by pulling back the idempotent matrix \eqref{eqn:universalidempotent} 
\[
M(P)=\begin{pmatrix} 1-x_1y_1 & -x_2y_1 & -x_3y_1 \\ -x_1y_2 & 1-x_2y_2 & -x_3y_2 \\ -x_1y_3 & -x_2y_3 & 1-x_3y_3\end{pmatrix}
\]
on $Q_5$ along $\mathbb{J}^3\xrightarrow{\pi_3} Q_6\xrightarrow{\Sigma_{\pone}\eta'}Q_5$.

\subsection{Invariants of the bundle}
We may compute the invariants associated to the above rank $2$ bundle by looking at its class in the set $[\mathbb{J}^3, \mathrm{BSL}_2]_{\A^1}=[\mathbb{P}^3, \mathrm{BSL}_2]_{\A^1}$.
Recall from \cite[\S 6]{Asok12a} that if $X$ is a smooth $k$-scheme of dimension $\leq 3$, there is an exact sequence of pointed sets
\[
\mathrm{H}^1(X,\piaone_2(\mathrm{BSL}_2))\to \mathrm{H}^3(X,\piaone_3(\mathrm{BSL}_2))\to [X,\mathrm{BSL}_2]_{\A^1}\to \mathrm{H}^2(X,\piaone_2(\mathrm{BSL}_2))\to *
\]
deduced from the Postnikov tower of $\mathrm{BSL}_2$ (using the fact that $\mathrm{BSL}_2$ is $\A^1$-simply connected). We know that $\piaone_2(\mathrm{BSL}_2)=\KMW{2}$, while the sheaf $\piaone_3(\mathrm{BSL}_2)$ fits into the exact sequence \eqref{eqn:pi2a2minus0}
\[
\KMW4/12_\epsilon\to \piaone_3(\mathrm{BSL}_2)\to \mathbf{GW}_3^2\to 0.
\]
We suppose that the base field $k$ is algebraically closed, in which case we obtain from \cite[Theorem 11.7]{Fasel09d}
\[
\mathrm{H}^3(\mathbb{P}^3,\KMW4/12_\epsilon)=\mathrm{H}^3(\mathbb{P}^3,\mathbf{K}^{\mathrm{M}}_4/12)=k^\times/(k^\times)^{12}=0,
\] 
and the above exact sequence then reads 
\[
\mathrm{H}^1(\mathbb{P}^3,\KMW2)\to \mathrm{H}^3(\mathbb{P}^3,\mathbf{GW}_3^2)\to [\mathbb{P}^3,\mathrm{BSL}_2]_{\A^1}\to \mathrm{H}^2(\mathbb{P}^3,\KMW2)\to *
\]
The inclusion $i\colon\mathbb{P}^3_k\smallsetminus \{0\}\subset \mathbb{P}^3_k$ (for which an affine model is given by the inclusion $\mathbb{J}^3\smallsetminus\{\A^3\}\to \mathbb{J}^3$) yields a commutative diagram
\[
\xymatrix@C=1em{\mathrm{H}^1(\mathbb{P}^3_k,\KMW2)\ar[r]\ar[d]^-{i^*} & \mathrm{H}^3(\mathbb{P}^3_k,\mathbf{GW}_3^2)\ar[r]\ar[d]^-{i^*} & [\mathbb{P}^3_k,\mathrm{BSL}_2]_{\A^1}\ar[r]\ar[d]^-{i^*} & \mathrm{H}^2(\mathbb{P}^3_k,\KMW2)\ar[r]\ar[d]^-{i^*} & \star \\
\mathrm{H}^1(\mathbb{P}^3_k\smallsetminus \{0\},\KMW2)\ar[r] & \mathrm{H}^3(\mathbb{P}^3_k\smallsetminus \{0\},\mathbf{GW}_3^2)\ar[r] & [\mathbb{P}^3_k\smallsetminus \{0\},\mathrm{BSL}_2]_{\A^1}\ar[r] & \mathrm{H}^2(\mathbb{P}^3_k\smallsetminus \{0\},\KMW2)\ar[r] & \star }
\]
The projection $\mathbb{P}^3\smallsetminus\{0\}\to \mathbb{P}^2$ defined by $[x_0:x_1:x_2:x_3]\mapsto [x_0:x_1:x_2]$ is a weak equivalence and consequently $\mathbb{P}^3\smallsetminus\{0\}$ is of $\A^1$-cohomological dimension $\leq 2$. The bottom sequence then reduces to a bijection
\[
[\mathbb{P}^3_k\smallsetminus \{0\},\mathrm{BSL}_2]_{\A^1}\to \mathrm{H}^2(\mathbb{P}^3_k\smallsetminus \{0\},\KMW2)
\] 
given by the Euler class. It follows from \cite[Theorem 11.7]{Fasel09d} that the projection map
\[
\mathrm{H}^2(\mathbb{P}^3,\KMW2)\to \mathrm{H}^2(\mathbb{P}^3,\mathbf{K}^{\mathrm{M}}_2)=\mathrm{CH^2}(\mathbb{P}^3)=\Z
\]
is an isomorphism. The same applies for $\mathbb{P}^2_k$ and thus the above diagram is of the form
\[
\xymatrix@C=1.5em{\mathrm{H}^1(\mathbb{P}^3_k,\KMW2)\ar[r] & \mathrm{H}^3(\mathbb{P}^3_k,\mathbf{GW}_3^2)\ar[r]\ar[d] & [\mathbb{P}^3_k,\mathrm{BSL}_2]_{\A^1}\ar[r]^-{c_2}\ar[d]^-{i^*} & \Z\ar[r]\ar@{=}[d] & \star \\
 & \star\ar[r] & [\mathbb{P}^2_k,\mathrm{BSL}_2]_{\A^1}\ar[r]_-{c_2} & \Z\ar[r] & \star }
\]
By construction, the bundle constructed using the exotic Hopf map is classified by the composite 
\[
\mathbb{J}^3\to Q_6\xrightarrow{\Sigma_{\pone}\eta'}Q_5\to \mathrm{BSL}_2
\]
The sequence $\mathbb{J}^3\smallsetminus \A^3\xrightarrow{i} \mathbb{J}^3\to Q_6$ being isomorphic to the cofiber sequence 
\[
\mathbb{P}^3_k\smallsetminus\{0\}\to \mathbb{P}^3_k\to (\pone)^{\wedge 3}
\]
we obtain that the composite $\mathbb{J}^3\smallsetminus \A^3\xrightarrow{i} \mathbb{J}^3\to Q_6$ is trivial up to homotopy. The diagram then shows in particular that the second Chern class of this bundle vanishes. We can now use the morphism $\mathbb{P}^3_k\to (\pone)^{\wedge 3}$, modeled by the morphism $\mathbb{J}^3\to Q_6$, to obtain a commutative diagram with exact columns
\[
\xymatrix@C=1.5em{\mathrm{H}^1((\pone)^{\wedge 3},\KMW2)\ar[r] & \mathrm{H}^3((\pone)^{\wedge 3},\mathbf{GW}_3^2)\ar[r]\ar[d] & [(\pone)^{\wedge 3},\mathrm{BSL}_2]_{\A^1}\ar[r]\ar[d]^-{i^*} & \mathrm{H}^2((\pone)^{\wedge 3},\KMW2)\ar[r]\ar@{=}[d] & \star \\
\mathrm{H}^1(\mathbb{P}^3_k,\KMW2)\ar[r] & \mathrm{H}^3(\mathbb{P}^3_k,\mathbf{GW}_3^2)\ar[r]\ar[d] & [\mathbb{P}^3_k,\mathrm{BSL}_2]_{\A^1}\ar[r]^-{c_2}\ar[d]^-{i^*} & \Z\ar[r]\ar@{=}[d] & \star \\
 & \star\ar[r] & [\mathbb{P}^2_k,\mathrm{BSL}_2]_{\A^1}\ar[r]_-{c_2} & \Z\ar[r] & \star }
\]
In view of \cite[Proposition 1.1.5]{Asok16}, the only nontrivial group in the top row is 
\[
\mathrm{H}^3((\pone)^{\wedge 3},\mathbf{GW}_3^2)=\mathbf{GW}_0^3(k)=\Z/2\Z,
\] and we thus obtain 
\[
\xymatrix@C=1.5em{\star\ar[r] & \mathbf{GW}_0^3(k)\ar[r]\ar[d] & [(\pone)^{\wedge 3},\mathrm{BSL}_2]_{\A^1}\ar[r]\ar[d]^-{i^*} & \star\ar[r]\ar@{=}[d] & \star \\
\mathrm{H}^1(\mathbb{P}^3_k,\KMW2)\ar[r] & \mathrm{H}^3(\mathbb{P}^3_k,\mathbf{GW}_3^2)\ar[r]\ar[d] & [\mathbb{P}^3_k,\mathrm{BSL}_2]_{\A^1}\ar[r]^-{c_2}\ar[d]^-{i^*} & \Z\ar[r]\ar@{=}[d] & \star \\
 & \star\ar[r] & [\mathbb{P}^2_k,\mathrm{BSL}_2]_{\A^1}\ar[r]_-{c_2} & \Z\ar[r] & \star }
\]
We claim that he composite
\[
\Z/2=\mathbf{GW}_0^3(k)=\mathrm{H}^3((\pone)^{\wedge 3},\mathbf{GW}_3^2)\to \mathrm{H}^3(\mathbb{P}^3_k,\mathbf{GW}_3^2)\to [\mathbb{P}^3_k,\mathrm{BSL}_2]_{\A^1}
\]
is split injective.

Indeed, we consider the map $[\mathbb{P}^3_k,\mathrm{BSL}_2]_{\A^1}\to [\mathbb{P}^3_k,\mathrm{BSp}]_{\A^1}=\widetilde{\mathrm{GW}}^2(\mathbb{P}_k^3)$ induced by the stabilization, where the right-hand side is the reduced Grothendieck-Witt group of symplectic bundles. The push-forward map induced by the projection $\mathbb{P}^3_k\to \Spec(k)$ reads as
\[
\widetilde{\mathrm{GW}}^2(\mathbb{P}_k^3)\to \mathrm{GW}^3(\Spec(k))=\Z/2
\]
and is explicitly described in \cite[\S 4.2.2]{Asok25a}, where it is called the motivic Atiyah-Rees invariant. The computation therein shows that this homomorphism provides a splitting of $\Z/2\to [\mathbb{P}^3_k,\mathrm{BSL}_2]_{\A^1}$, and thus that the Atiyah-Rees invariant of the exotic Hopf bundle is $1\in \Z/2$.


%%%%%%%%%%%%%%%%%%%%%%%%%%%%%%%%%%%%%%%%%%%%%%%%%%%%%%%%%%%%%%%%%%%%%%%%%%%%%%%%%%%%%%%%

\begin{footnotesize}
\bibliographystyle{alpha}
\bibliography{General}%,OtherReferences}
\end{footnotesize}
\Addresses
\end{document}
